% Options for packages loaded elsewhere
\PassOptionsToPackage{unicode}{hyperref}
\PassOptionsToPackage{hyphens}{url}
\PassOptionsToPackage{dvipsnames,svgnames,x11names}{xcolor}
\documentclass[
  11pt,
]{article}
\usepackage{xcolor}
\usepackage[margin=2.5cm]{geometry}
\usepackage{amsmath,amssymb}
\setcounter{secnumdepth}{-\maxdimen} % remove section numbering
\usepackage{iftex}
\ifPDFTeX
  \usepackage[T1]{fontenc}
  \usepackage[utf8]{inputenc}
  \usepackage{textcomp} % provide euro and other symbols
\else % if luatex or xetex
  \usepackage{unicode-math} % this also loads fontspec
  \defaultfontfeatures{Scale=MatchLowercase}
  \defaultfontfeatures[\rmfamily]{Ligatures=TeX,Scale=1}
\fi
\usepackage{lmodern}
\ifPDFTeX\else
  % xetex/luatex font selection
\fi
% Use upquote if available, for straight quotes in verbatim environments
\IfFileExists{upquote.sty}{\usepackage{upquote}}{}
\IfFileExists{microtype.sty}{% use microtype if available
  \usepackage[]{microtype}
  \UseMicrotypeSet[protrusion]{basicmath} % disable protrusion for tt fonts
}{}
\makeatletter
\@ifundefined{KOMAClassName}{% if non-KOMA class
  \IfFileExists{parskip.sty}{%
    \usepackage{parskip}
  }{% else
    \setlength{\parindent}{0pt}
    \setlength{\parskip}{6pt plus 2pt minus 1pt}}
}{% if KOMA class
  \KOMAoptions{parskip=half}}
\makeatother
\usepackage{longtable,booktabs,array}
\newcounter{none} % for unnumbered tables
\usepackage{calc} % for calculating minipage widths
% Correct order of tables after \paragraph or \subparagraph
\usepackage{etoolbox}
\makeatletter
\patchcmd\longtable{\par}{\if@noskipsec\mbox{}\fi\par}{}{}
\makeatother
% Allow footnotes in longtable head/foot
\IfFileExists{footnotehyper.sty}{\usepackage{footnotehyper}}{\usepackage{footnote}}
\makesavenoteenv{longtable}
\setlength{\emergencystretch}{3em} % prevent overfull lines
\providecommand{\tightlist}{%
  \setlength{\itemsep}{0pt}\setlength{\parskip}{0pt}}
\usepackage[utf8]{inputenc}
\usepackage[T1]{fontenc}
\usepackage{amsmath}
\usepackage{amssymb}
\usepackage{textcomp}
\usepackage{newunicodechar}

% --- Greek letters ---
\newunicodechar{Δ}{\ensuremath{\Delta}}
\newunicodechar{Σ}{\ensuremath{\Sigma}}
\newunicodechar{α}{\ensuremath{\alpha}}
\newunicodechar{β}{\ensuremath{\beta}}
\newunicodechar{δ}{\ensuremath{\delta}}
\newunicodechar{η}{\ensuremath{\eta}}
\newunicodechar{θ}{\ensuremath{\theta}}
\newunicodechar{ι}{\ensuremath{\iota}}
\newunicodechar{μ}{\ensuremath{\mu}}
\newunicodechar{π}{\ensuremath{\pi}}
\newunicodechar{ρ}{\ensuremath{\rho}}
\newunicodechar{σ}{\ensuremath{\sigma}}
\newunicodechar{τ}{\ensuremath{\tau}}
\newunicodechar{φ}{\ensuremath{\varphi}}
\newunicodechar{χ}{\ensuremath{\chi}}
\newunicodechar{ψ}{\ensuremath{\psi}}

% --- Math operators ---
\newunicodechar{⟨}{\ensuremath{\langle}}
\newunicodechar{⟩}{\ensuremath{\rangle}}
\newunicodechar{∈}{\ensuremath{\in}}
\newunicodechar{∉}{\ensuremath{\notin}}
\newunicodechar{⊥}{\ensuremath{\perp}}
\newunicodechar{↦}{\ensuremath{\mapsto}}
\newunicodechar{∐}{\ensuremath{\coprod}}
\newunicodechar{≠}{\ensuremath{\neq}}
\newunicodechar{≅}{\ensuremath{\cong}}
\newunicodechar{⊕}{\ensuremath{\oplus}}
\newunicodechar{⊗}{\ensuremath{\otimes}}
\newunicodechar{⊆}{\ensuremath{\subseteq}}
\newunicodechar{∀}{\ensuremath{\forall}}
\newunicodechar{∃}{\ensuremath{\exists}}
\newunicodechar{∅}{\ensuremath{\emptyset}}
\newunicodechar{⟺}{\ensuremath{\Longleftrightarrow}}
\newunicodechar{−}{\ensuremath{-}}
\newunicodechar{∧}{\ensuremath{\wedge}}
\newunicodechar{∪}{\ensuremath{\cup}}
\newunicodechar{≈}{\ensuremath{\approx}}
\newunicodechar{≡}{\ensuremath{\equiv}}
\newunicodechar{≥}{\ensuremath{\geq}}
\newunicodechar{↑}{\ensuremath{\uparrow}}
\newunicodechar{→}{\ensuremath{\rightarrow}}
\newunicodechar{↓}{\ensuremath{\downarrow}}
\newunicodechar{↔}{\ensuremath{\leftrightarrow}}
\newunicodechar{′}{\ensuremath{\prime}}
\newunicodechar{×}{\ensuremath{\times}}
\newunicodechar{¬}{\ensuremath{\neg}}

% --- Subscripts ---
\newunicodechar{₀}{\textsubscript{0}}
\newunicodechar{₁}{\textsubscript{1}}
\newunicodechar{₂}{\textsubscript{2}}
\newunicodechar{₃}{\textsubscript{3}}
\newunicodechar{₄}{\textsubscript{4}}
\newunicodechar{₅}{\textsubscript{5}}
\newunicodechar{₆}{\textsubscript{6}}
\newunicodechar{₇}{\textsubscript{7}}
\newunicodechar{ₖ}{\textsubscript{k}}

% --- Status indicators ---
\newunicodechar{✅}{\textbf{[YES]}}
\newunicodechar{⚠}{\textbf{[WARN]}}
\newunicodechar{❌}{\textbf{[NO]}}

% --- Sanskrit / Pali diacritics (not covered by T1 inputenc) ---
\newunicodechar{ḥ}{\d{h}}
\newunicodechar{ṃ}{\d{m}}
\newunicodechar{ṅ}{\.{n}}
\newunicodechar{ṇ}{\d{n}}
\newunicodechar{ṣ}{\d{s}}

% --- Combining characters ---
\newunicodechar{̂}{\^{}}

% --- Miscellaneous ---
\newunicodechar{–}{--}
\newunicodechar{—}{---}
\newunicodechar{·}{\textperiodcentered}
\newunicodechar{²}{\textsuperscript{2}}
\newunicodechar{°}{\textdegree}
\newunicodechar{§}{\S}
\newunicodechar{─}{---}
\newunicodechar{️}{}
\usepackage{bookmark}
\IfFileExists{xurl.sty}{\usepackage{xurl}}{} % add URL line breaks if available
\urlstyle{same}
\hypersetup{
  colorlinks=true,
  linkcolor={Maroon},
  filecolor={Maroon},
  citecolor={Blue},
  urlcolor={blue},
  pdfcreator={LaTeX via pandoc}}

\author{}
\date{}

\begin{document}

Author: VietVunVut (Viet - Nguyen Xuan); GitHub:
https://github.com/AIhugART/; Facebook:
https://www.facebook.com/xuanviet

\section{When Does a Physical Interaction Become a Valid Registered
Measurement?}\label{when-does-a-physical-interaction-become-a-valid-registered-measurement}

\subsection{A VVV-QMRF Registration-Layer Framework with the K9\_E Class
C Testable Hypothesis and an Experimental Specification for Extended
Wigner's
Friend}\label{a-vvv-qmrf-registration-layer-framework-with-the-k9_e-class-c-testable-hypothesis-and-an-experimental-specification-for-extended-wigners-friend}

\textbf{Working Paper v3.0} \emph{(promoted from draft 2026-05-28 ---
all phases P1-P7 complete)} \textbf{Author:} Viet Nguyen Xuan
(VietVunVut) \textbf{Affiliation:} Independent Researcher, Vietnam
\textbf{Contact:} viet@vvvqmrf.com \textbar{} https://vvvqmrf.com
\textbf{Repository:}
https://github.com/AIhugART/VVV-QMRF-VietVunVut-Quantum-Measurement-Registration-Framework
\textbf{Date:} 2026-05-28 \textbf{Status:} Working paper. All formal
claims are Class D (proposed) or Class C (conjecture/qualified) unless
stated otherwise. Critique is explicitly invited. \textbf{Plan
reference:}
\texttt{papers/paper\_003/VVV-QMRF\_Working\_Paper\_v3.0\_plan.md}
\textbf{Base version:} v2.0
(\texttt{papers/Testable\_Prediction\_Section/.../VVV-QMRF\_Working\_Paper\_v2.0.md})

\begin{quote}
\textbf{DISCLAIMER:} VVV-QMRF is independent personal research, not
Standard Quantum Mechanics, not peer-reviewed, and not experimentally
confirmed. K9\_E (P9) is classified \textbf{Class C (qualified)}:
structurally testable but empirically unconfirmed. Evidence is real but
ambiguous; noise as an alternative explanation cannot be ruled out (v30
noise sensitivity analysis FAIL). Confirmation or rejection requires a
purpose-designed K9-S12 photonic Extended Wigner's Friend experiment.
VVV-QMRF does not replace Standard QM, revise the Born rule, or invoke
consciousness. Full boundary protocol: \texttt{DISCLAIMER.md}. Formal
definitions:
\texttt{documents/research\_documents/project\_vvv\_qmrf\_class\_c/06\_references/VVV\_QMRF\_Definitions.md}.
\end{quote}

\begin{center}\rule{0.5\linewidth}{0.5pt}\end{center}

\subsection{Abstract}\label{abstract}

Standard quantum mechanics specifies, via postulate P3, that a
measurement of observable A on state \textbar ψ⟩ yields eigenvalue aₖ
with probability \textbar⟨aₖ\textbar ψ⟩\textbar². What P3 does not
specify is \emph{when} a physical interaction constitutes a valid
registered measurement event. This silence generates the von Neumann
chain problem --- every measuring apparatus becomes entangled with the
system it measures, with no postulate indicating where this chain
terminates --- and leaves the Heisenberg cut formally undefined.

This paper proposes VVV-QMRF (VietVunVut Quantum Measurement
Registration Framework): a registration-layer extension of quantum
mechanics grounded in Buddhist Pramāṇa epistemology
(Dignāga--Dharmakīrti tradition). The framework introduces a K-side
registration space K, structurally separate from the physical Hilbert
space H (K ≠ H), and proposes a six-condition test for valid registered
measurement. These six conditions formalize when a physical interaction
crosses from the ρ-side into the K-side as a self-certified,
validity-bearing registration event. Applied to Extended Wigner's Friend
(EWF) scenarios, this yields the K-side incommensurability conjecture:
K\_F ⊥\_K K\_W holds when a joint registration space K\_joint is
structurally required but cannot satisfy all six conditions
simultaneously for both observers.

The K-side registration space is axiomatized via eight frozen
registration-logic axioms (K1--K8, Layer 1) and nine updatable bridge
theorems (T1--T9, Layer 2). K1 defines the K-state carrier set
(five-tuple: M, o, cert, t, V); K2 imposes a strict total temporal
order; K3 formalizes self-certification; K4 assigns default validity; K5
governs invalidation and cross-registration contradiction; K6 defines
cross-registration authority; K7 formalizes registration process closure
(provisional → final validity); K8 guarantees validity preservation
under cross-space embeddings. Layer 2 extends this structure with bridge
theorems T1--T9, including K5\_prospective (a conservative
pre-instantiation evaluation mode required for probability assignment),
K7\_trace (a closure transition record), and D\_enc (a
transition-encoding semantic definition). T4-H --- the N-observer
colimit existence hypothesis --- has been verified as a full theorem
across all four steps (RCA 4.74/5, 2026-05-28), upgrading the
three-observer prediction to Class C.

K9\_E is the probability postulate (Postulate P9) motivated by the
K-space architecture. It is not a theorem derivable from K1--K8 alone
--- the core axioms define structural properties (registration,
validity, incommensurability) but do not uniquely determine a
probability rule. K9\_E fills this gap as P9: P(o\textbar K) = Tr(E\_o
ρ) · (1 − β·K\_ctx), where β ∈ {[}0,1{]} is a single free parameter and
K\_ctx is the fraction of contextual observer registrations in
prospective incommensurability with the current registration. At β = 0,
K9\_E reduces exactly to the Born rule. Six of the eight terms in K9\_E
are entirely new concepts not present in Standard QM. The proposed
experimental specification (K9-S12 Modified Bong Protocol --- single
waveplate at angle α = 31°) predicts Genuine LF 1 = +0.0891 (8.6σ),
δ⟨A₁B₂⟩ = −0.0355 (20.8σ), and figure of merit FOM = 8.6; this proposal
was submitted to arXiv as a separate experimental child paper
(paper\_002, 2026-05-27).

\textbf{Empirical status (Class C qualified):} A genuine non-circular
fit to raw Proietti et al.~(2019) data yields β = 0.598, visibility V =
0.939, Δχ² = 5.35 (2.31σ) in favor of K9\_E over quantum mechanics with
uniform visibility. However, a v30 noise sensitivity analysis returned
FAIL: the noise threshold for evidence robustness is 0.10 σ RMS (well
below the 3.0 σ threshold for PASS). Random noise at any magnitude
produces Δχ² ≥ 5.35 in approximately 50\% of realizations due to K9\_E's
directional sensitivity and the small sample of four data points. The
A0B0 setting alone drives 80\% of the Δχ² improvement. The 2.31σ signal
reflects model flexibility, not confirmed physical suppression.
Adversarial tests (4/4 PASS) and operationalizability gates (3/3 PASS)
confirm K9\_E's structural robustness, but empirical confirmation
remains open and requires the K9-S12 photonic EWF experiment with
dedicated noise characterization.

This paper presents three logically independent projects under a one-way
motivation chain. \textbf{Project A} (BE↔QM Comparative Mapping)
provides a 30-node, 39-edge comparative framework between Buddhist
Pramāṇa epistemology and quantum measurement --- an interpretive
framework that motivated the registration-layer architecture but does
not itself constitute evidence for that architecture. \textbf{Project B}
(VVV-QMRF Conceptual Framework) develops the formal K-space
axiomatization (K1--K8, Layer 1 frozen) and bridge theorem suite
(T1--T9, Layer 2 updatable), including the conjectured
structure-preserving map φ: K → B(H) with necessary conditions N₁--N\_T
(Class D supporting conjecture). \textbf{Project C} (K9\_E Testable
Hypothesis) is the scientifically falsifiable claim: the K9\_E
probability postulate (P9) produces predictions structurally different
from Standard QM (δ\_S ≠ 0 when β \textgreater{} 0), avoids the
Frauchiger--Renner paradox via K5 V\_prov mechanism, and is falsifiable
via the K9-S12 Modified Bong Protocol. A→B→C is a one-way motivation
chain, not a logical derivation. A null K9\_E result falsifies Project C
but does not invalidate Projects A or B.

VVV-QMRF does not replace Standard Quantum Mechanics. It does not revise
postulates P1--P4 or the Born rule (K9\_E recovers Born exactly at β =
0). It does not invoke consciousness. It does not claim Buddhist
epistemology proves quantum mechanics. All formal claims are classified
by evidence level; classification criteria are stated in the project's
formal definitions document. Falsification requires a photonic EWF
experiment with the K9-S12 single-waveplate protocol or equivalent.
Critique is explicitly invited.

\begin{center}\rule{0.5\linewidth}{0.5pt}\end{center}

\subsection{1. The Registration Layer
Gap}\label{the-registration-layer-gap}

\subsubsection{1.1 What Quantum Mechanics
Specifies}\label{what-quantum-mechanics-specifies}

Standard quantum mechanics is built on four postulates:

\begin{itemize}
\tightlist
\item
  \textbf{P1 (State):} A physical system is represented by a density
  operator ρ ∈ D(H).
\item
  \textbf{P2 (Observables):} Physical quantities are represented by
  self-adjoint operators A on H.
\item
  \textbf{P3 (Measurement):} Measurement of A on state \textbar ψ⟩
  yields eigenvalue aₖ with probability \textbar⟨aₖ\textbar ψ⟩\textbar²,
  after which the state updates to the corresponding eigenstate.
\item
  \textbf{P4 (Dynamics):} Between measurements, the state evolves
  unitarily via the Schrödinger equation.
\end{itemize}

P3 specifies what outcome appears and with what probability. P3 does not
specify when a physical interaction counts as a measurement, what
structural conditions distinguish a measurement from a non-measurement
interaction, or what stops the measurement chain. This is the
registration layer gap.

\subsubsection{1.2 The Von Neumann Chain
Problem}\label{the-von-neumann-chain-problem}

Von Neumann (1932) observed that if apparatus A1 measures system S, A1
becomes entangled with S. If apparatus A2 then measures A1, A2 becomes
entangled with A1+S. No postulate of standard QM specifies where this
chain terminates. Decoherence (Zurek 2003) explains why superpositions
become effectively classical at the macroscopic scale, but it does not
explain when a physical interaction constitutes a registration event in
the sense of a definite recorded outcome. The von Neumann chain problem
is a registration-layer problem, not a decoherence problem.

\subsubsection{1.3 The Heisenberg Cut}\label{the-heisenberg-cut}

Copenhagen quantum mechanics (Bohr 1935; Heisenberg 1958) assigns
measurement authority to ``the classical apparatus'' without formally
defining what makes an apparatus classical, why it has registration
authority, or what structural conditions produce a definite outcome.
VVV-QMRF addresses both problems by introducing a registration layer K,
separate from the physical layer ρ, and proposing formal conditions for
valid registration events within K. The K1--K8 axioms (§4) formalize
this registration layer structurally.

\begin{center}\rule{0.5\linewidth}{0.5pt}\end{center}

\subsection{2. The K-side Registration
Space}\label{the-k-side-registration-space}

\subsubsection{2.1 Layer Separation}\label{layer-separation}

VVV-QMRF introduces a strict separation between two layers:

\begin{itemize}
\tightlist
\item
  \textbf{Physical layer (ρ-side):} Physical states, observables, and
  dynamics as described by Standard QM. This layer is not modified.
\item
  \textbf{Registration layer (K-side):} The space of registration
  events, recording conditions, and validity states. K ≠ H. K-side
  structure does not alter ρ-side dynamics.
\end{itemize}

\subsubsection{2.2 Minimal K-state}\label{minimal-k-state}

A registration event is represented by a minimal K-state tuple:

\begin{verbatim}
k = ⟨M, o, cert, t, V⟩
\end{verbatim}

The K-side registration space K is the collection of such tuples
produced by a registering system R over time. Section §4 formalizes this
minimal K-state as the K1--K8 axiom set.

\subsubsection{2.3 Source: Buddhist Pramāṇa Epistemology (Project
A)}\label{source-buddhist-pramux101ux1e47a-epistemology-project-a}

The registration-layer architecture is structurally derived from
Buddhist Pramāṇa epistemology (Dignāga, 5th century; Dharmakīrti, 7th
century). The structural extraction concerns: \textbf{Svasaṃvedana}
(self-certifying cognition, source for K3), \textbf{Arthakriyā}
(validity as functional success, source for K4/K7), and
\textbf{Anumāna/Pratyakṣa} (inference and perception as distinct
registration modes). This is Project A: a comparative interpretive
framework. Project A motivates Project B via one-way structural
inspiration; it does not constitute empirical evidence for the K-space
architecture or K9\_E.

\begin{center}\rule{0.5\linewidth}{0.5pt}\end{center}

\subsection{3. The Valid Registered Measurement
Test}\label{the-valid-registered-measurement-test}

\subsubsection{3.1 The Six Conditions}\label{the-six-conditions}

\textbf{Interaction X is a valid registered measurement for registering
system R if and only if:}

\begin{verbatim}
Condition 1 (Physical):    X occurs at the physical ρ-side.
Condition 2 (Admission):   X is admitted into K-side as M_X for system R.
Condition 3 (Process):     M_X ∈ R ordered by registration time.          [K2, E6]
Condition 4 (Self-cert):   σ_R(M_X) = 1 intrinsically within K_R.         [K3, E1]
Condition 5 (Validity):    V(M_X) = 1 by default upon admission.           [K4, E7]
Condition 6 (Non-invalid): No later M′ revises V(M_X) to 0.               [K5, K7, E7]
\end{verbatim}

The cert and V fields are formalized as K3 (self-certification) and
K4/K5/K7 (default validity and invalidation) in §4.

\subsubsection{3.2 The K-side Stopping
Proposition}\label{the-k-side-stopping-proposition}

\textbf{Proposition (Class D proposed):} If Conditions 1--6 hold for
M\_X, no further K-side registration act is required to certify M\_X.
The von Neumann registration regress terminates at M\_X on the K-side
without modifying ρ-side dynamics. Condition 4 (self-certification)
provides the formal property that Copenhagen assigns to ``classical
apparatus'' without definition.

\subsubsection{3.3 Observer-Indexed
Self-Certification}\label{observer-indexed-self-certification}

For multi-observer scenarios: σ\_R(M) = 1 iff M occurred as a K-side
registration event of R, determined intrinsically within K\_R, not by
any M′ ≠ M and not by any R′ ≠ R. The independence σ\_F(M\_F) ⊥
σ\_W(M\_W) is formalized in K3 (§4.1).

\begin{center}\rule{0.5\linewidth}{0.5pt}\end{center}

\subsection{4. K-Space Architecture: Axioms K1--K8 and Layer
2}\label{k-space-architecture-axioms-k1k8-and-layer-2}

This section presents the formal axiomatization of the K-side
registration space. The axioms and bridge theorems are the structural
backbone of VVV-QMRF and the foundation on which K9\_E (§5) and the EWF
incommensurability result (§6) rest.

\begin{quote}
\textbf{Scope boundary (v2.4 §0.6 audit):} The K-space axiomatization in
this section is \textbf{purely structural}. It contains zero probability
equations, zero numerical values, and zero experimental data. K9\_E (§5)
is a separate probability postulate (P9) that is \textbf{not derivable
from K1--K8 alone} --- the axioms define structural properties
(registration, validity, incommensurability) but do not uniquely
determine a probability rule. Data fitting (§7) is a separate empirical
analysis. This three-layer separation is architectural: K1--K8
(structure) → K9\_E (postulate) → data fitting (empirical).
\end{quote}

The authoritative peer-synced source is
\texttt{documents/research\_documents/meta\_architecture/K\_Space\_Axiomatization.md}
(canonical) and
\texttt{documents/research\_documents/project\_vvv\_qmrf\_class\_c/01\_axiomatization/K\_Space\_Axiomatization.md}
(working copy), both v2.4.

\subsubsection{4.1 Layer 1 (Frozen): Axioms
K1--K8}\label{layer-1-frozen-axioms-k1k8}

Layer 1 contains eight registration-logic axioms whose text is frozen
unconditionally. K1--K4 and K8 carry no Level 4 semantic dependencies.
K5, K6, and K7 have conditional semantic dependencies on Level 4 scope
(when they fire / when closure occurs) but their text is frozen
regardless of Level 4 changes.

{\def\LTcaptype{none} % do not increment counter
\begin{longtable}[]{@{}
  >{\raggedright\arraybackslash}p{(\linewidth - 8\tabcolsep) * \real{0.2000}}
  >{\raggedright\arraybackslash}p{(\linewidth - 8\tabcolsep) * \real{0.2000}}
  >{\raggedright\arraybackslash}p{(\linewidth - 8\tabcolsep) * \real{0.2000}}
  >{\raggedright\arraybackslash}p{(\linewidth - 8\tabcolsep) * \real{0.2000}}
  >{\raggedright\arraybackslash}p{(\linewidth - 8\tabcolsep) * \real{0.2000}}@{}}
\toprule\noalign{}
\begin{minipage}[b]{\linewidth}\raggedright
Axiom
\end{minipage} & \begin{minipage}[b]{\linewidth}\raggedright
Name
\end{minipage} & \begin{minipage}[b]{\linewidth}\raggedright
Core function
\end{minipage} & \begin{minipage}[b]{\linewidth}\raggedright
BE lineage
\end{minipage} & \begin{minipage}[b]{\linewidth}\raggedright
Class
\end{minipage} \\
\midrule\noalign{}
\endhead
\bottomrule\noalign{}
\endlastfoot
\textbf{K1} & Carrier Set & K\_R = set of 5-tuples k = ⟨M, o, cert, t,
V⟩; cert = 1 for all k ∈ K\_R (admission invariant); t-injectivity:
distinct events have distinct timestamps & Pramāṇa --- cognition as
structured event: act, object, self-awareness, result, validity & C \\
\textbf{K2} & Temporal Order & (K\_R, \textless\_R) is strict total
order (chain): k₁ \textless\_R k₂ iff t(k₁) \textless{} t(k₂); discrete:
no K-side registration identity between consecutive events &
Kṣaṇabhaṅgavāda --- momentariness; registration time is discrete within
K\_R & D \\
\textbf{K3} & Self-Certification & cert(k) = σ\_R(M) ∈ \{0,1\}
determined intrinsically within K\_R; independent of any R′ ≠ R;
observer-indexed independence & Svasaṃvedana --- self-certifying
cognition: a cognition certifies its own occurrence without a
second-order cognition & D \\
\textbf{K4} & Default Validity & V(k) = 1 upon instantiation for
¬isNull(k); V(k) = 0 for null events (isNull: o = ∅ ∧ ΔI = 0); V is
provisional until K7 closure & Svataḥ prāmāṇya --- intrinsic validity by
default; arthakriyā (causal efficacy) & D \\
\textbf{K5} & Invalidation & V(k₁) → 0 iff ∃k₂ later in K2 order, k₂ ⊥
k₁ in shared C\_K, k₂ has Auth(K6); fires only when requires\_K\_joint =
1; pre-closure reversible, post-closure absolute & Parataḥ prāmāṇya ---
invalidity detected extrinsically; bādhaka pramāṇa & D \\
\textbf{K6} & Cross-Reg. Authority & Auth(k₂ → k₁, C\_K) = 1 iff: both
in same C\_K; V(k₂) = 1; k₁ ∈ scope(D\_joint); non-transitive across
distinct C\_K & Bādhaka pramāṇa --- a contradicting cognition must
itself be a valid pramāṇa & D \\
\textbf{K7} & Closure & R closes at t\_close when no pending
requires\_K\_joint demands remain; V\_prov → V\_final; post-closure: no
new k, K5 irreversibility absolute, no new D\_joint & Niścaya
(ascertainment) --- a cognition becomes determinate at cognitive closure
& D \\
\textbf{K8} & Embedding Preservation & For embedding i: K\_R → K\_X:
V\_X(i(k)) = V\_R(k) at t\_embed; all five fields preserved;
post-embedding V evolves under K\_X's K4--K7 & Anugama --- continuity:
validity accompanies a cognition into a broader context & D \\
\end{longtable}
}

\textbf{Note on K1 claim class:} K1 is classified C because the K-state
tuple definition is shared with the K9\_E postulate (which is Class C);
K1 provides K9\_E's domain. K2--K8 remain Class D as structural axioms
without empirical content.

\subsubsection{4.2 K5\_prospective: Conservative Pre-Instantiation
Evaluation
Extension}\label{k5_prospective-conservative-pre-instantiation-evaluation-extension}

K9\_E requires assessing incommensurability \emph{before} an outcome is
registered --- to compute the suppression factor f\_perp(o) = P(K5 fires
prospectively on hypothetical outcome o). K5's standard post-hoc mode
(modifying V of an actual tuple) is not designed for this.
K5\_prospective is a \textbf{conservative extension} that adds a new
evaluation target while preserving K5's identical conditions (i)--(iii):

\begin{verbatim}
K5 (post-hoc):     V(k₁) → 0 iff ∃k₂: k₁ <_R k₂ ∧ k₂ ⊥ k₁ ∧ Auth(k₂→k₁)
                   Target: actual k₁ ∈ K_R.    Effect: V_prov modified.

K5_prospective:    K5 fires on k_o* iff ∃k_prev: k_prev <_joint k_o*
                     ∧ k_o* ⊥ k_prev ∧ Auth(k_o*→k_prev)
                   Target: hypothetical k_o* = ⟨M*, o, cert=1, t*, V=1⟩.
                   Effect: contributes to f_perp(o) in K9_E only.
\end{verbatim}

\textbf{Same conditions (i)--(iii). Same structural logic. Different
target and effect.} K5\_prospective does NOT modify V of any actual
tuple in K\_R. It has no operational role outside K9\_E. This
conservative extension was upgraded from ``semantic extension'' to an
explicit axiom-level clause in v29 (2026-05-23, 3-Round RCA Round 2
score 4.90/5), resolving assumption {[}A-E1{]} through the T9
construction theorem.

\textbf{Claim class:} C (conservative extension of K5; same conditions;
new evaluation mode only).

\subsubsection{4.3 Layer 2 (Updatable): Bridge Theorems
T1--T9}\label{layer-2-updatable-bridge-theorems-t1t9}

Layer 2 bridge theorems connect K1--K8 to the framework's operational
definitions (Level 4 predicates: D\_joint, requires\_K\_joint,
Bridge\_EWF, ODC\_K). Layer 2 is updatable independently of K1--K8.

{\def\LTcaptype{none} % do not increment counter
\begin{longtable}[]{@{}
  >{\raggedright\arraybackslash}p{(\linewidth - 8\tabcolsep) * \real{0.2000}}
  >{\raggedright\arraybackslash}p{(\linewidth - 8\tabcolsep) * \real{0.2000}}
  >{\raggedright\arraybackslash}p{(\linewidth - 8\tabcolsep) * \real{0.2000}}
  >{\raggedright\arraybackslash}p{(\linewidth - 8\tabcolsep) * \real{0.2000}}
  >{\raggedright\arraybackslash}p{(\linewidth - 8\tabcolsep) * \real{0.2000}}@{}}
\toprule\noalign{}
\begin{minipage}[b]{\linewidth}\raggedright
Theorem
\end{minipage} & \begin{minipage}[b]{\linewidth}\raggedright
Name
\end{minipage} & \begin{minipage}[b]{\linewidth}\raggedright
Core function
\end{minipage} & \begin{minipage}[b]{\linewidth}\raggedright
Class
\end{minipage} & \begin{minipage}[b]{\linewidth}\raggedright
Status
\end{minipage} \\
\midrule\noalign{}
\endhead
\bottomrule\noalign{}
\endlastfoot
\textbf{T1} & K\_joint Construction & K\_joint exists as categorical
colimit of embedding diagram when requires\_K\_joint = 1; order =
transitive closure of embedded orders + cross-temporal relations & D &
Pending L4 freeze \\
\textbf{T2} & ⊥\_K Derivation & K\_A ⊥\_K K\_B iff requires\_K\_joint =
1 ∧ no admissible K\_joint; traced to K5 conflict (V\_prov forced to 0
in candidate K\_joint) or K7 lock & D & Pending L4 freeze \\
\textbf{T3} & Bridge\_EWF & Bridge\_EWF = 1 derivable from K5 in EWF:
M\_F registers definite o\_F; M\_W registers same lab as superposition;
no reinterpretation preserves both without changing validity claim;
conditional on AJVS & D/C & Pending L4 freeze \\
\textbf{T4} & N-Observer Colimit & K\_joint(R₁,\ldots,R\_N) exists as
colimit for N ≥ 2 when pairwise AdmJoint + global commutativity hold;
⊥\_K non-transitive in general & D & \textbf{T4-H THEOREM 4/4
(2026-05-28)} \\
\textbf{T5} & K\_joint Associativity & K\_joint(K\_joint(A,B),C) ≅
K\_joint(A,B,C) as K1--K8-preserving isomorphism; conditional on T4-H &
D & T4-H resolved; L4 freeze pending \\
\textbf{T6} & Decoherence Update & Decoherence-induced ρ-side change
produces K-side registration-state update when Conditions 1--6 hold & D
& Updatable \\
\textbf{T7} & IRB Scope Propagation & Intrinsic Relational Binding
(entanglement) between R\_i and R\_j induces requires\_K\_joint = 1
under specific conditions & D & Updatable \\
\textbf{T8} & K5\_prospective Frequency Bridge & f\_perp =
E{[}I(K5\_prospective fires){]} = fraction of K\_ctx with prospective ⊥;
derives K9\_E suppression term structurally; eliminates assumption
{[}A-E2{]} & C & Structural derivation from K5p \\
\textbf{T9} & K\_ctx Construction & K\_ctx = \{φ\_ij(k\_j) ∈ K\_joint :
k\_j ∈ K\_\{R\_j\}, requires\_K\_joint = 1, temporally compatible\};
φ\_ij = canonical K8-constrained T1 embedding; eliminates {[}A-E1{]} via
5 lemmas (L1-L5) & C & L1--L5 proven \\
\end{longtable}
}

\textbf{T4-H THEOREM (4/4, 2026-05-28, RCA 4.74/5):} All four steps
verified: - Step 1: C\_\{K-space\} satisfies category axioms (identity,
composition, associativity). - Step 2: K\_colim = (∐\_i
K\_i)/\textasciitilde{} constructed; all five fields well-defined via
lexicographic t-assignment and embedding-time V snapshot; 5/5
verification gates PASS. - Step 3: K1--K8 preserved through quotient
construction via T-PRES Lemma + T-REP Corollary (K2 acyclicity resolved)
+ V monotone dynamics. - Step 4: Universal property verified:
u({[}k,i{]}) := f\_i(k) is well-defined (diagram compatibility + T-REP),
K8-preserving, and unique.

T4-H is no longer a hypothesis. T4 conclusions hold for all N ≥ 2. The
three-observer prediction (§10) is upgraded from Class C-conditional to
Class C. T1 (N=2 constructive) remains independently valid without
invoking the colimit universal property.

\subsubsection{4.4 K7\_trace and D\_enc: Canonical Layer 2 (v2.4,
2026-05-27)}\label{k7_trace-and-d_enc-canonical-layer-2-v2.4-2026-05-27}

\textbf{K7\_trace --- Closure Transition Record:} A conservative
extension of K7. Defines Δ\_closure(k) := V\_prov(k) − V\_final(k) at
t\_close. Records the transition magnitude (0 or 1) without modifying V
or creating new tuples. Claim class: C-canonical. Promoted from BB-VVV
local origin to canonical Layer 2 by RCA 4.77/5.

\textbf{D\_enc --- Transition-Encoding Registration Act:} A semantic
definition (parent: K7\_trace). Enc(M\_aware, k\_F) = 1 iff o(M\_aware)
depends counterfactually on Δ\_closure(k\_F). Pattern: same as
K5\_prospective (binary classification of hypothetical act). Claim
class: C-canonical.

\textbf{Four canonical consumers of K7\_trace (cross-validation of
architectural status):}

{\def\LTcaptype{none} % do not increment counter
\begin{longtable}[]{@{}
  >{\raggedright\arraybackslash}p{(\linewidth - 4\tabcolsep) * \real{0.3333}}
  >{\raggedright\arraybackslash}p{(\linewidth - 4\tabcolsep) * \real{0.3333}}
  >{\raggedright\arraybackslash}p{(\linewidth - 4\tabcolsep) * \real{0.3333}}@{}}
\toprule\noalign{}
\begin{minipage}[b]{\linewidth}\raggedright
Consumer
\end{minipage} & \begin{minipage}[b]{\linewidth}\raggedright
Usage
\end{minipage} & \begin{minipage}[b]{\linewidth}\raggedright
Status
\end{minipage} \\
\midrule\noalign{}
\endhead
\bottomrule\noalign{}
\endlastfoot
T\_BB V3 Step 1 (BB-VVV) & Δ\_closure as V\_prov substitute post-closure
(Baumann-Brukner 2024) & §9.1 --- T\_BB Class C conditional \\
D\_enc & Parent for Enc predicate & Canonical Layer 2 \\
3-OBS hierarchical registration transition & T4-H THEOREM unlocks Class
C upgrade for 3-observer prediction & §10 --- Class C (2026-05-28) \\
FR-VVV avoidance chain Step 2 & Frauchiger-Renner K5→K6→V=0 path;
K7\_trace scenario-agnostic across B\&B and FR & §9.2 --- V\_FR2 PASS
(2026-05-28) \\
\end{longtable}
}

The four-consumer canonical status confirms K7\_trace is an
architectural element, not a single-use construct.

\subsubsection{4.5 K-Space and Buddhist Pramāṇa: One-Way Motivation, Not
Empirical
Evidence}\label{k-space-and-buddhist-pramux101ux1e47a-one-way-motivation-not-empirical-evidence}

The BE lineage entries in §4.1 record the structural inspiration for
each axiom. This is Project A's role in VVV-QMRF: a one-way motivation
from Pramāṇa theory's analysis of valid cognition to K-space's analysis
of valid registration.

\textbf{What this does NOT mean:} - Buddhist epistemology does not prove
K1--K8. - K1--K8 are not derived from Buddhist texts; they are derived
from the registration-layer problem posed by Standard QM. - Project A's
30-node, 39-edge BE↔QM comparative mapping is a separate independent
project (see §12.3 for the A→B→C independence statement). - A result
that falsifies K9\_E does not falsify Buddhist epistemology, the BE↔QM
mapping, or the K-space architecture.

\begin{center}\rule{0.5\linewidth}{0.5pt}\end{center}

\subsection{5. K9\_E: Probability Postulate
(P9)}\label{k9_e-probability-postulate-p9}

\begin{quote}
\textbf{ERRATUM note (from Phase 8 source document, 2026-05-23):} This
section was formerly titled ``K9\_E Formal Derivation.'' K9\_E is
\textbf{NOT derived from K1--K8}. It is a \textbf{POSTULATE} --- a
probability assignment rule motivated by K-space structure (⊥\_K,
K\_ctx) but not uniquely determined by K1--K8 axioms. This distinction
is load-bearing: the entire §5 must be read as proposing P9, not proving
it.
\end{quote}

\subsubsection{5.1 Why K1--K8 Alone Do Not Determine a Probability
Rule}\label{why-k1k8-alone-do-not-determine-a-probability-rule}

The eight axioms of Layer 1 define the structural properties of the
K-side registration space: what tuples are admitted (K1), how they are
ordered (K2), how self-certification works (K3), how default validity is
assigned (K4), how cross-registration contradiction fires (K5), what
counts as authority for invalidation (K6), when the registration process
closes (K7), and how validity is preserved under embedding (K8).

Nowhere in K1--K8 is there a formula that maps K-state tuples to
real-valued probabilities. The axioms determine:

\begin{itemize}
\tightlist
\item
  \emph{Which} registration events are valid (cert = 1, V = 1)
\item
  \emph{When} one event invalidates another (K5 conditions)
\item
  \emph{How} events embed across K-spaces (K8)
\end{itemize}

They do not determine \emph{how much probability} to assign to each
outcome. Multiple distinct probability functions are consistent with the
K1--K8 structure. In particular, the Born rule P(o) = Tr(E\_o ρ) is
consistent with K1--K8 (with β = 0, the K-space constraints are
satisfied trivially). But so is any function that respects the V-gate (V
= 0 → no probability), the isNull-gate (isNull(k) → no probability), and
normalizes over the outcome space.

K9\_E fills this gap by adding a specific probability assignment rule as
\textbf{Postulate P9}, motivated by the K-space incommensurability
structure (K5\_prospective, T8, T9) but not derivable from it. This is
analogous to how the Born rule P3 in Standard QM is an irreducible
postulate of the theory --- it is not derivable from P1 (state), P2
(observables), or P4 (dynamics).

\subsubsection{5.2 The K9\_E Equation}\label{the-k9_e-equation}

\textbf{Postulate P9 (K9\_E Probability Postulate):}

\begin{verbatim}
P(o | k_i, Exp) = Tr(E_o ρ_i) · [1 − β · f_perp(o, k_i, K_ctx)]
                  ─────────────────────────────────────────────────
                                    Z_E(k_i)

where:
  f_perp(o, k_i, K_ctx) = |{k_j ∈ K_ctx : K5_prospective fires on k_o* vs k_j}|
                           ────────────────────────────────────────────────────────
                                              |K_ctx|

  K_ctx(k_i, Exp) = {φ_ij(k_j) ∈ K_joint : k_j ∈ K_{R_j},
                     requires_K_joint(R_i, R_j) = 1,
                     temporally compatible with k_i}        [T9]

  Z_E(k_i) = Σ_{o'} Tr(E_{o'} ρ_i) · [1 − β · f_perp(o')]  [normalization]

  β ∈ [0, 1)  — single free parameter (suppression strength)
\end{verbatim}

\textbf{Born rule limit (exact, not approximate):}

K9\_E reduces to the Born rule exactly under any of the following
equivalent conditions:

\begin{verbatim}
P(o|k) = Tr(E_o ρ)   iff any of:

  (i)   β = 0                  [suppression OFF — K-space has no effect on P]
  (ii)  K_ctx = ∅              [no contextual observers — f_perp undefined → 0 by convention]
  (iii) f_perp(o) = 0 ∀ o     [all outcomes compatible — no prospective K5 firing]
  (iv)  Single observer N=1   [implies K_ctx = ∅ → (ii)]
\end{verbatim}

For condition (i): β = 0 → {[}1 − 0·f\_perp{]} = 1 for all o → Z\_E = Σ
Tr(E\_o ρ) = 1 → P = Tr(E\_o ρ). Standard single-observer laboratory
experiments have N = 1 (condition iv) and therefore K9\_E is
observationally equivalent to Standard QM in all non-EWF contexts.

\subsubsection{5.3 Eight-Term Provenance: What Is New in
K9\_E}\label{eight-term-provenance-what-is-new-in-k9_e}

K9\_E assembles eight structural components. Six of these eight terms
are entirely new concepts not present in Standard QM. This table records
the provenance of each term --- it is a traceability record, not a
derivation proof.

{\def\LTcaptype{none} % do not increment counter
\begin{longtable}[]{@{}
  >{\raggedright\arraybackslash}p{(\linewidth - 8\tabcolsep) * \real{0.2000}}
  >{\raggedright\arraybackslash}p{(\linewidth - 8\tabcolsep) * \real{0.2000}}
  >{\raggedright\arraybackslash}p{(\linewidth - 8\tabcolsep) * \real{0.2000}}
  >{\raggedright\arraybackslash}p{(\linewidth - 8\tabcolsep) * \real{0.2000}}
  >{\raggedright\arraybackslash}p{(\linewidth - 8\tabcolsep) * \real{0.2000}}@{}}
\toprule\noalign{}
\begin{minipage}[b]{\linewidth}\raggedright
\#
\end{minipage} & \begin{minipage}[b]{\linewidth}\raggedright
Term
\end{minipage} & \begin{minipage}[b]{\linewidth}\raggedright
Definition
\end{minipage} & \begin{minipage}[b]{\linewidth}\raggedright
Source
\end{minipage} & \begin{minipage}[b]{\linewidth}\raggedright
In Standard QM?
\end{minipage} \\
\midrule\noalign{}
\endhead
\bottomrule\noalign{}
\endlastfoot
T1 & \texttt{Tr(E\_o\ ρ\_i)} & Born rule probability (POVM formulation)
& \textbf{Standard QM} & ✅ \\
T2 & \texttt{β} & Suppression strength, β ∈ {[}0,1) & \textbf{Free
parameter} --- the single adjustable parameter of K9\_E; analogous to
coupling constants (α, G, g) in physics & ❌ \textbf{NEW} \\
T3 & \texttt{f\_perp(o,\ k\_i,\ K\_ctx)} & Fraction of contextual
registrations prospectively incommensurable with outcome o &
\textbf{K5\_prospective + T8} --- structural derivation from binary
K5/K6 primitives (T8-H1: fraction form is uniquely forced) & ❌
\textbf{NEW} \\
T4 & \texttt{C(o\_i,\ o\_j)} & Compatibility map: 1 if outcomes are
QM-orthogonal (incompatible), 0 otherwise & \textbf{K-space
construction} from ρ\_joint at setup, stored as K-side lookup & ❌
\textbf{NEW} \\
T5 & \texttt{K\_ctx(k\_i,\ Exp)} & Set of contextual K-states from other
observers via T9 morphism + temporal compatibility & \textbf{T9 + K2}
--- K\_ctx is a theorem, not an assumption (L1--L5 in T9) & ❌
\textbf{NEW} \\
T6 & \texttt{Z\_E(k\_i)} & Normalization: Σ\_\{o'\} Tr(E\_\{o'\}
ρ\_i)·{[}1−β·f\_perp(o'){]} & Modified from QM (Standard QM
auto-normalizes; K9\_E suppression breaks auto-normalization, requiring
explicit Z) & ⚠️ \textbf{MODIFIED} \\
T7 & \texttt{V(k)\ =\ 0\ →\ no\ P} & Bhrānti gate: invalid registrations
receive no probability & \textbf{K4 + K5 → PP-1 v2} --- grounded in BE
concept bhrānti (erroneous cognition) & ❌ \textbf{NEW} \\
T8 & \texttt{isNull(k)\ →\ no\ P} & Anupalabdhi gate: null events (o =
∅, ΔI = 0) receive no probability & \textbf{K4 isNull guard} ---
grounded in BE concept anupalabdhi (non-apprehension) & ❌
\textbf{NEW} \\
\end{longtable}
}

\textbf{Summary:} Of eight terms in K9\_E: 1 from Standard QM (T1: Born
rule), 1 modified from QM (T6: normalization), \textbf{6 entirely new}
(T2--T5, T7, T8). The four original K9\_E construction assumptions
({[}A-E1{]}--{[}A-E4{]}) are fully resolved: {[}A-E1{]} ELIMINATED (T9,
5 lemmas), {[}A-E2{]} ELIMINATED (T8/T8-H1), {[}A-E3{]} RECLASSIFIED as
free parameter β, {[}A-E4{]} BE-anchored. \textbf{Net: 0 orphaned
assumptions, 1 free parameter (β).}

\subsubsection{5.4 K9\_E Is a Postulate, Not a
Theorem}\label{k9_e-is-a-postulate-not-a-theorem}

\textbf{This is the most important boundary statement in this section.}

K9\_E is \textbf{Postulate P9} --- a proposed probability assignment
rule added to the VVV-QMRF framework alongside K1--K8. It is not
derivable from K1--K8 alone. The axioms define structural properties
(registration, validity, incommensurability) but do not uniquely
determine a probability rule.

The relationship between K1--K8 and K9\_E is the same as the
relationship between P1, P2, P4 (state, observables, dynamics) and P3
(Born rule) in Standard QM: the structural postulates define the
mathematical space, but the probability rule is an additional postulate
with its own empirical motivation.

\textbf{What K1--K8 do:} Define the K-space structural constraints
(carrier, order, certification, validity, invalidation, authority,
closure, embedding). Supply the structural vocabulary used in K9\_E
(K5\_prospective, T8, T9). Constrain the form of a probability rule
(V-gate via T7, isNull-gate via T8), but do not determine β or the
specific suppression formula.

\textbf{What K9\_E adds:} A specific probability assignment rule
P(o\textbar K) with a single free parameter β. The suppression term
f\_perp is structurally motivated by K5\_prospective incommensurability
and formally derived (via T8, T9) to have the fraction form --- but the
\emph{presence} of the suppression term (i.e., that β ≠ 0 in nature) is
a postulate, not a derivation.

\textbf{Consequence:} K9\_E is falsifiable (β = 0 recovers Born exactly;
any measured β ≠ 0 confirms K9\_E; K9-S12 protocol provides the
definitive test). K1--K8 are not individually falsifiable --- they are
structural definitions of the registration layer.

\subsubsection{5.5 Distinguishability from Standard
QM}\label{distinguishability-from-standard-qm}

When β \textgreater{} 0, K9\_ctx ≠ ∅, and f\_perp is outcome-dependent,
K9\_E produces predictions structurally different from Standard QM.

\textbf{Distinguishability condition:} δP(o) = P\_\{K9E\}(o) − Tr(E\_o
ρ) ≠ 0 iff ALL of:

\begin{verbatim}
  (I)   β > 0
  (II)  K_ctx ≠ ∅   (requires_K_joint = 1 for at least one observer pair)
  (III) ∃ o, o' : f_perp(o) ≠ f_perp(o')
\end{verbatim}

Condition (III) is critical: if f\_perp is constant across outcomes, the
suppression factor {[}1 − β·f\_perp{]} is a uniform multiplier that
cancels in the normalized ratio P/Z\_E → Tr(E\_o ρ). This is the PP-2 v2
cancellation insight: K9\_E must suppress \emph{differentially} across
outcomes to produce an observable signal.

In EWF scenarios with Bell-State Measurement (BSM), the compatibility
map C(o\_F, o\_W) is outcome-dependent (some (o\_F, o\_W) pairs are
compatible, others are incommensurable), producing outcome-dependent
f\_perp and therefore a genuine δP ≠ 0.

\textbf{Theoretical magnitude (CHSH, β = 0.5):} δ\_S = −0.055. This is a
\emph{theoretical} value --- the predicted deviation in the CHSH
parameter for Proietti-type 2-observer EWF. At β = 0.9, δ\_S = −0.179 (≈
2.4σ detection horizon). No empirical detection of δ\_S has been
achieved to date. See §7 for empirical status.

\textbf{3-observer amplification (Class C, T4-H THEOREM):} The
3-observer K\_ctx is larger, producing an amplification factor of
approximately 11× relative to 2-observer. Predicted δ\_M3 = −0.223 at β
= 0.3 (illustrative, conditional only on K9\_E postulate P9 --- T4-H is
now THEOREM, §10).

\subsubsection{5.6 Projects B and C Are Independent
Claims}\label{projects-b-and-c-are-independent-claims}

The φ-map conjecture (§11, Project B) and K9\_E (Project C) are
logically independent:

\textbf{Project B (φ-map):} Conjectures the existence of a
structure-preserving map φ: K → B(H). Derives necessary conditions
N₁--N\_T. Provides φ-conditional scope boundaries for existing QM
interpretations. This conjecture is Class D and concerns the
\emph{formal relationship} between K-space and the algebra of bounded
operators. It does not make testable predictions about measurement
statistics.

\textbf{Project C (K9\_E):} Proposes a probability postulate P9 with a
single free parameter β. Makes testable predictions (K9-S12). Can be
confirmed or falsified by experiment. K9\_E's empirical status is
entirely independent of whether φ exists.

\textbf{Independence in both directions:} - A proof that φ: K → B(H)
exists (or does not exist) would not confirm or falsify K9\_E. - A
K9-S12 experimental confirmation of K9\_E (β ≠ 0 measured) would not
prove φ exists. - A K9-S12 null result (β = 0 measured) would falsify
K9\_E but would not falsify the φ-map conjecture or the K-space
architecture.

This independence is architectural: Project B concerns the structural
embedding of K into H; Project C concerns the probability assignment
within K. They share the K-space vocabulary (K1--K8, T1--T9) but their
claims are logically disjoint.

\begin{center}\rule{0.5\linewidth}{0.5pt}\end{center}

\subsection{6. Extended Wigner's Friend and K-side
Incommensurability}\label{extended-wigners-friend-and-k-side-incommensurability}

\subsubsection{6.1 The Extended Wigner's Friend
Setup}\label{the-extended-wigners-friend-setup}

In the Extended Wigner's Friend scenario (Frauchiger and Renner 2018;
Brukner 2018), two observer pairs interact with a shared quantum system:

\begin{itemize}
\tightlist
\item
  \textbf{Friend F} measures quantum system S inside a sealed
  laboratory. From F's perspective, a definite outcome o\_F is recorded:
  σ\_F(M\_F) = 1.
\item
  \textbf{Wigner W} models the entire laboratory (F + S) as a unitarily
  evolving quantum state, then performs an interference measurement M\_W
  on F's laboratory.
\end{itemize}

Standard QM describes both perspectives as valid without providing a
formal account of their structural incompatibility. This is where the
registration layer offers new structure.

\subsubsection{6.2 Applying the Six Conditions to Each
Observer}\label{applying-the-six-conditions-to-each-observer}

\textbf{Friend F:}

\begin{verbatim}
X_F → M_F ∈ R_F         [Condition 3, K2]
σ_F(M_F) = 1             [Condition 4, K3]
V_F(M_F) = 1             [Condition 5, K4]
No M′ contradicts M_F    [Condition 6, K5/K7]
→ M_F is a valid registered measurement within K_F.
\end{verbatim}

\textbf{Wigner W:}

\begin{verbatim}
X_W → M_W ∈ R_W         [Condition 3, K2]
σ_W(M_W) = 1             [Condition 4, K3]
V_W(M_W) = 1             [Condition 5, K4]
No M′ contradicts M_W    [Condition 6, K5/K7]
→ M_W is a valid registered measurement within K_W.
\end{verbatim}

Both observers satisfy the six conditions independently within their own
K-sides. The question is whether a joint registration space K\_joint can
contain both K\_F and K\_W as jointly valid entries.

\subsubsection{6.3 The requires\_K\_joint
Predicate}\label{the-requires_k_joint-predicate}

\textbf{Definition (admissible joint K-side registration space, Class D
proposed):}

\begin{verbatim}
AdmJoint(K_joint; A, B) = 1
  iff  there exist embeddings i_A: A → K_joint and i_B: B → K_joint such that:
    (i)   embeddings preserve act, outcome, cert, registration time/order, validity;
    (ii)  self-certification remains intrinsic to each embedded act;
    (iii) Conditions 1-6 remain satisfied for each embedded structure;
    (iv)  no required registration-state update in K_joint invalidates either embedded
          structure while both are still claimed as jointly valid;
    (v)   K_joint does not import an external certifier as source of self-certification.
\end{verbatim}

\textbf{Definition (requires\_K\_joint predicate, Class D proposed):}

\begin{verbatim}
requires_K_joint(A, B) = 1
  iff  A and B are each valid or provisionally valid within their own K-side
  AND  A and B are brought under a shared validity demand D_joint
  AND  D_joint requires both to be assessed as parts of the same registration target,
       history, counterfactual claim, or validity claim
  AND  truth of D_joint cannot be evaluated while leaving A and B in fully independent K-sides
  AND  preserving D_joint requires a K_joint in which A and B are jointly valid.

requires_K_joint(A, B) = 0
  iff  no shared D_joint is imposed, or D_joint can be evaluated without embedding A and B
       into one candidate K_joint.
\end{verbatim}

\textbf{Operational sufficient conditions for requires\_K\_joint = 1:}

\begin{itemize}
\tightlist
\item
  \textbf{Condition A (Wigner interference):} W performs an interference
  measurement on F's lab; M\_W registers a superposition description
  while M\_F registers a definite outcome for the same lab. →
  requires\_K\_joint = 1.
\item
  \textbf{Condition B (Direct comparison):} F and W directly compare
  records and a logical contradiction is detectable. →
  requires\_K\_joint = 1.
\item
  \textbf{Condition B2 (LF constraint):} An LF inequality requires
  F-side and W-side claims to have simultaneous cross-observer validity.
  → requires\_K\_joint = 1.
\end{itemize}

\textbf{Operational sufficient conditions for requires\_K\_joint = 0:}

\begin{itemize}
\tightlist
\item
  \textbf{Condition C (No interference, no comparison):} W does not
  perform interference on F's lab. → requires\_K\_joint = 0.
\item
  \textbf{Condition D (Separable state):} Shared state is separable;
  M\_F and M\_W act on non-overlapping subsystems. → requires\_K\_joint
  = 0.
\item
  \textbf{Condition E (Independent bookkeeping):} Records stored
  together but no joint validity demand imposed. → requires\_K\_joint =
  0.
\end{itemize}

\subsubsection{6.4 Formal Definition of K-side
Incommensurability}\label{formal-definition-of-k-side-incommensurability}

**Definition (K-side incommensurability relation ⊥\_K, Class D
proposed):**

\begin{verbatim}
A ⊥_K B
  iff  requires_K_joint(A, B) = 1
  AND  there exists no admissible K_joint such that:
       (i)   A and B both embed into K_joint,
       (ii)  their respective self-certifications are preserved,
       (iii) their respective validity conditions remain satisfied, and
       (iv)  no required registration-state update in K_joint invalidates
             either A or B while both are still claimed as jointly valid.
\end{verbatim}

\textbf{Definition (K-side comparison context C\_K):} A minimal shared
frame in which two registration acts are evaluated for compatibility.
C\_K requires: (a) both acts admitted into the same comparison domain;
(b) both indexed to the same registration target; (c) comparison does
not presuppose joint validity. C\_K is strictly weaker than AdmJoint.

\textbf{Definition (cross-registration authority):} In C\_K, M\_later
has valid cross-registration authority w.r.t. M\_earlier iff: (a) both
belong to same C\_K; (b) V(M\_later) = 1; (c) M\_later's content
directly concerns the same registration target as M\_earlier; (d) the
comparison architecture does not arbitrarily privilege one observer.

\textbf{Boundary clauses:} - ⊥\_K does not assert that either physical
event fails to occur on the ρ-side. - ⊥\_K does not mean either
observer's outcome is false within its own K-side. - ⊥\_K is not
equivalent to a registration-null event Null\_K(e). - ⊥\_K applies only
when both sides are valid/provisionally valid within their own K-sides.

\subsubsection{6.5 K\_joint Failure and K-side
Incommensurability}\label{k_joint-failure-and-k-side-incommensurability}

\textbf{Bridge lemma (Class D proposed; application status: Class C):}

\begin{verbatim}
Bridge_EWF(D_joint; M_F, M_W) = 1
  iff  D_joint requires F-side and W-side registrations to be evaluated as jointly valid
       parts of one laboratory registration history
  AND  M_F registers a definite friend-side outcome o_F
  AND  M_W registers the same laboratory as coherent superposition (no definite o_F as W-side claim)
  AND  LF/no-go comparison requires both claims to support one cross-observer validity constraint
  AND  no reinterpretation inside K_joint preserves both without changing the validity claim
       of at least one side.

Bridge_EWF(D_joint; M_F, M_W) = 1 → M_W ⊥ M_F.
\end{verbatim}

\textbf{Relativization defense (rejected):} A relativized K\_joint
hosting only meta-descriptions (``within K\_F, M\_F registered
\textbar h⟩'') does not satisfy D\_joint. D\_joint requires joint
validity of the \emph{original} registration claims, not
meta-descriptions of them. Relativizing the contents abandons D\_joint
rather than satisfying it.

\textbf{Conditional lemma (K\_joint failure, Class D proposed):} Under
requires\_K\_joint = 1, Bridge\_EWF = 1, and valid cross-registration
authority, no admissible K\_joint exists such that σ\_F(M\_F) = 1,
σ\_W(M\_W) = 1, V(M\_F) = 1, and V(M\_W) = 1 simultaneously hold. K\_F
⊥\_K K\_W follows. \textbf{Claim class of Step 4:} C/D boundary ---
conditional on full Bridge\_EWF semantic proof.

\subsubsection{6.6 The Falsifiable Prediction and
ODC\_K}\label{the-falsifiable-prediction-and-odc_k}

VVV-QMRF conjectures that configurations satisfying requires\_K\_joint =
1 are the natural candidates for LF-level violation when the quantum
state is sufficiently entangled.

\textbf{Operational data criterion (ODC\_K, Class C proposed):}

\begin{verbatim}
ODC_K(Data, Cfg) = K_joint_fails
  iff  Cfg satisfies requires_K_joint = 1 via D_joint and Bridge_EWF,
  AND  no joint registration model J_K satisfying:
       (i)   jointly valid F-side and W-side registrations,
       (ii)  Conditions 1-3, 5, 6 preserved (Condition 4 by construction),
       (iii) AdmJoint condition (iv): no required invalidation while both are jointly claimed,
       (iv)  observed probability distribution reproduced within predeclared tolerance τ,
       (v)   not reclassifying a Bridge_EWF config as mere independent bookkeeping.
\end{verbatim}

τ must be fixed before data collection; ODC\_K is sensitive to the
choice of τ.

\textbf{Compatibility checks with existing data:}

\emph{Proietti et al.~(2019):} The three correlators with
requires\_K\_joint = 1 (⟨A₁B₁⟩, ⟨A₁B₀⟩, ⟨A₀B₁⟩) contribute positively to
the violated expression S = 2.407; the one term with requires\_K\_joint
= 0 (⟨A₀B₀⟩) is subtracted. This allocation is structurally compatible.
The VVV-QMRF reading is a term-role interpretation of an aggregate
inequality, not a per-configuration experimental confirmation.

\emph{Bong et al.~(2020):} Two-regime structure:

{\def\LTcaptype{none} % do not increment counter
\begin{longtable}[]{@{}
  >{\raggedright\arraybackslash}p{(\linewidth - 4\tabcolsep) * \real{0.3333}}
  >{\raggedright\arraybackslash}p{(\linewidth - 4\tabcolsep) * \real{0.3333}}
  >{\raggedright\arraybackslash}p{(\linewidth - 4\tabcolsep) * \real{0.3333}}@{}}
\toprule\noalign{}
\begin{minipage}[b]{\linewidth}\raggedright
μ regime
\end{minipage} & \begin{minipage}[b]{\linewidth}\raggedright
Reported result
\end{minipage} & \begin{minipage}[b]{\linewidth}\raggedright
VVV-QMRF reading
\end{minipage} \\
\midrule\noalign{}
\endhead
\bottomrule\noalign{}
\endlastfoot
μ = 0.80-0.81 & Bell non-LF violated; LF inequalities NOT violated &
Regime 1: requires\_K\_joint = 1 active, entanglement insufficient for
LF failure \\
μ ≈ 0.87 & First LF violation (Semi-Brukner) & Transition Regime 1 →
2 \\
High μ & All categories violated including Genuine LF & Regime 2:
K\_joint failure exposed \\
\end{longtable}
}

\textbf{Per-facet μ-threshold ODC\_K prediction:}

{\def\LTcaptype{none} % do not increment counter
\begin{longtable}[]{@{}
  >{\raggedright\arraybackslash}p{(\linewidth - 4\tabcolsep) * \real{0.3333}}
  >{\raggedright\arraybackslash}p{(\linewidth - 4\tabcolsep) * \real{0.3333}}
  >{\raggedright\arraybackslash}p{(\linewidth - 4\tabcolsep) * \real{0.3333}}@{}}
\toprule\noalign{}
\begin{minipage}[b]{\linewidth}\raggedright
Facet class
\end{minipage} & \begin{minipage}[b]{\linewidth}\raggedright
μ=0.80-0.81
\end{minipage} & \begin{minipage}[b]{\linewidth}\raggedright
High-μ
\end{minipage} \\
\midrule\noalign{}
\endhead
\bottomrule\noalign{}
\endlastfoot
Bell non-LF & K\_joint\_fails (violated) & K\_joint\_fails (violated) \\
Brukner & K\_joint\_exists (not violated) & K\_joint\_fails
(violated) \\
Semi-Brukner & K\_joint\_exists (not violated) & K\_joint\_fails
(violated) \\
Genuine LF & K\_joint\_exists (not violated) & K\_joint\_fails
(violated) \\
\end{longtable}
}

The prediction hierarchy (Bell non-LF first, then Semi-Brukner, then
Genuine LF) matches the observed μ-threshold ordering. \textbf{Not yet a
facet-level confirmation} --- full ODC\_K stage 3 model-fit test has not
been performed. These are compatibility checks.

\subsubsection{6.7 Connection to K1--K8 and
K5\_prospective}\label{connection-to-k1k8-and-k5_prospective}

The Level 4 formalism of §6.1--6.6 can now be traced directly to Layer 1
axioms via the K-space architecture (§4):

\textbf{requires\_K\_joint = K5 precondition:}

\begin{verbatim}
requires_K_joint(A, B) = 1
  ⟺  K5 firing precondition met: comparison context C_K exists
  ⟺  D_joint imposes a shared validity demand on K_F and K_W
     (K5 formal block: "fires only when requires_K_joint = 1")
\end{verbatim}

K5 does not fire in isolation --- it fires only when there is a joint
validity demand. The requires\_K\_joint = 0/1 distinction in §6 is
precisely the K5 precondition distinction in Layer 1. Every operational
sufficient condition A--E for requires\_K\_joint maps to whether K5's
precondition is met.

\textbf{Bridge\_EWF = K5\_prospective evaluation:}

\begin{verbatim}
Bridge_EWF(D_joint; M_F, M_W) = 1
  ⟺  K5_prospective fires on hypothetical k_o* vs k_prev = k_F:
     — k_F <_joint k_o*         [K2 temporal order in K_joint from T1]
     — k_o* ⊥ k_F within C_K   [K5 registered contradiction, minimal ⊥ definition]
     — Auth(k_o* → k_F, C_K) = 1  [K6 authority]
\end{verbatim}

Bridge\_EWF = 1 means exactly that K5\_prospective would fire if W's
hypothetical registration k\_o* were realized --- the structural bridge
between Level 4 formalism and Layer 1 axioms. This is the connection
exploited by T8 to derive f\_perp (§5.3).

\textbf{K7\_trace and K\_joint failure:} When K\_F ⊥\_K K\_W is
established, K7\_trace records Δ\_closure(k\_F) := V\_prov(k\_F) −
V\_final(k\_F). The four canonical consumers of K7\_trace (§4.4, §9.3)
all derive from this same closure event in structurally distinct EWF
scenarios.

\begin{center}\rule{0.5\linewidth}{0.5pt}\end{center}

\subsection{7. Empirical Status (Class C
Qualified)}\label{empirical-status-class-c-qualified}

\textbf{Classification:} K9\_E = \textbf{Class C (qualified)} ---
structurally testable, empirically \textbf{UNCONFIRMED}. Evidence is
real (a genuine non-circular fit shows 2.31σ improvement over
QM-uniform-visibility) but ambiguous (noise at any magnitude produces an
equivalent improvement in \textasciitilde50\% of realizations).
Confirmation or rejection requires the K9-S12 purpose-designed photonic
EWF experiment.

\subsubsection{7.1 D1 Proietti: Genuine Non-Circular
Fit}\label{d1-proietti-genuine-non-circular-fit}

\textbf{Version history (critical):} The Phase 10 data fitting document
(2026-05-23) identified a \emph{circular fit} error in the earlier Phase
10 analysis: the ``data'' used were reconstructed as E\_exp = V\_exp ·
E\_QM, mathematically guaranteeing β = 0 as the best-fit result (a
tautology, not an empirical finding). The v29 analysis (2026-05-23)
replaced this with a \textbf{genuine non-circular fit} using raw
correlator values extracted directly from Proietti et al.~(2019) Figure
3 (\texttt{Wigner\_figure\_3.md} SOT document).

\textbf{Genuine fit results (raw Proietti Figure 3 data):}

{\def\LTcaptype{none} % do not increment counter
\begin{longtable}[]{@{}
  >{\raggedright\arraybackslash}p{(\linewidth - 4\tabcolsep) * \real{0.3333}}
  >{\raggedright\arraybackslash}p{(\linewidth - 4\tabcolsep) * \real{0.3333}}
  >{\raggedright\arraybackslash}p{(\linewidth - 4\tabcolsep) * \real{0.3333}}@{}}
\toprule\noalign{}
\begin{minipage}[b]{\linewidth}\raggedright
Quantity
\end{minipage} & \begin{minipage}[b]{\linewidth}\raggedright
Value
\end{minipage} & \begin{minipage}[b]{\linewidth}\raggedright
Meaning
\end{minipage} \\
\midrule\noalign{}
\endhead
\bottomrule\noalign{}
\endlastfoot
β (best-fit) & \textbf{0.598} & K9\_E suppression parameter at genuine
fit \\
V (visibility, fitted) & \textbf{0.939} & Higher than circular-fit V =
0.854; non-uniform visibility detected \\
χ²/DOF & 0.670 (DOF = 2) & Good fit quality, p = 0.51 \\
Δχ² (K9\_E vs QM-uniform-V) & \textbf{5.35} & K9\_E improves fit over QM
with uniform visibility \\
σ-significance & \textbf{2.31σ} & Marginal detection level \\
\end{longtable}
}

\textbf{Why the raw data differs from reconstructed data:} The raw
Proietti Figure 3 correlator values (e.g., A0B0 = −0.678) differ
substantially from reconstructed values (−0.604 in the circular-fit
analysis). The non-uniform pattern across the four correlators is what
drives the genuine-fit improvement --- when K9\_E's directional
suppression is applied, it preferentially reduces some correlators
(those with high requires\_K\_joint = 1 contribution) over others
(requires\_K\_joint = 0), producing a better fit than a
uniform-visibility model.

\textbf{Two K9\_E model variants co-exist} with different calibration
paths: - \textbf{Additive model} (\texttt{utils/k9e\_predictor.py}): E =
E\_QM · {[}1 − β·n\_BSM·g\_ctx{]}, g\_ctx ≈ 0.039 (from theoretical δ\_S
= −0.055 at β = 0.5). - \textbf{Multiplicative model}
(\texttt{proietti\_raw\_fit.py}): E = E\_QM · {[}1 −
β·g\_eff{]}\^{}n\_BSM, g\_eff = 0.146 (from PP-4 sanity check). Produces
larger suppression; used for genuine fit.

The two models agree at first order in β·g but diverge for β
\textgreater{} 0.3.

\subsubsection{7.2 v30 Noise Sensitivity Analysis:
FAIL}\label{v30-noise-sensitivity-analysis-fail}

The v30 noise sensitivity analysis (P10-NOISE, 2026-05-24) asked:
\emph{How large must non-uniform noise be to produce Δχ² ≥ 5.35 by
chance?}

\textbf{Result:} \textbf{FAIL.} The noise threshold is 0.10 σ RMS ---
far below the 3.0 σ threshold required for PASS.

{\def\LTcaptype{none} % do not increment counter
\begin{longtable}[]{@{}
  >{\raggedright\arraybackslash}p{(\linewidth - 4\tabcolsep) * \real{0.3333}}
  >{\raggedright\arraybackslash}p{(\linewidth - 4\tabcolsep) * \real{0.3333}}
  >{\raggedright\arraybackslash}p{(\linewidth - 4\tabcolsep) * \real{0.3333}}@{}}
\toprule\noalign{}
\begin{minipage}[b]{\linewidth}\raggedright
Metric
\end{minipage} & \begin{minipage}[b]{\linewidth}\raggedright
Value
\end{minipage} & \begin{minipage}[b]{\linewidth}\raggedright
Interpretation
\end{minipage} \\
\midrule\noalign{}
\endhead
\bottomrule\noalign{}
\endlastfoot
Noise threshold (2σ, B4 Monte Carlo) & \textbf{0.10 σ RMS} & Noise at
ANY magnitude produces Δχ² ≥ 5.35 in \textasciitilde50\% of
realizations \\
PASS threshold & 3.0 σ RMS & Threshold for ``noise ruled out'' \\
Verdict & \textbf{FAIL} (0.10 \textless\textless{} 3.0) & Non-uniform
noise cannot be ruled out as explanation \\
Single-setting fragility (B2) & \textbf{1.85 σ at A0B0} & Only a 1.85σ
shift at A0B0 eliminates K9\_E advantage \\
A0B0 contribution to Δχ² & \textbf{80\%} & Nearly the entire K9\_E
``signal'' is driven by one data point \\
\end{longtable}
}

\textbf{What this means:} The 2.31σ signal reflects \textbf{model
flexibility} (K9\_E has one more degree of freedom than
QM-uniform-visibility), not confirmed physical K9\_E suppression. With
only 4 correlator values and K9\_E's directional sensitivity, random
noise at any magnitude produces a Δχ² ≥ 5.35 improvement approximately
half the time. The genuine fit improvement is real but not
distinguishable from a noise artifact with current data.

\textbf{Why this downgraded K9\_E from ``genuine'' to ``qualified'':}
The v29 upgrade to Class C (genuine) was based on three conditions: (1)
genuine non-circular fit, (2) K5\_prospective axiom upgrade, (3) T4-H
Step 1. Conditions (2) and (3) remain valid. Condition (1) still holds
as a genuine fit, but the noise analysis reveals that the fit evidence
is insufficient to rule out noise. The v30 P10-NOISE FAIL downgraded the
empirical sub-condition from ``real evidence'' to ``real but ambiguous
evidence.'' Class C structural qualification is unchanged; the (genuine)
→ (qualified) downgrade concerns the empirical confidence level only.

\textbf{Methodology rationale (conservative):} The noise analysis was
designed to ERR ON THE SIDE OF CAUTION --- a Type I error (falsely
claiming noise is ruled out) is fatal to project credibility; a Type II
error (correctly noting that noise cannot be ruled out) is recoverable
via the K9-S12 experiment.

\subsubsection{7.3 K9E-PAT: CLOSED as
UNRESOLVABLE}\label{k9e-pat-closed-as-unresolvable}

\textbf{Open item K9E-PAT} asked whether the K9\_E multiplicative
suppression pattern (predicted 2BSM/1BSM ratio ≈ 2) is confirmed by
Proietti data.

\textbf{Result:} \textbf{CLOSED as UNRESOLVABLE (v31 RCA 4.92/5).}

The empirical 2BSM/1BSM ratio of −0.78 is not a meaningful test of the
K9\_E pattern:

\begin{enumerate}
\def\labelenumi{\arabic{enumi}.}
\tightlist
\item
  The ratio −0.78 is computed from A0B1 + A1B0 avg = −0.0235 divided by
  A1B1 = +0.0179 --- both are sub-σ residuals (consistent with zero).
\item
  The ratio of two near-zero numbers carries no physical information
  about suppression structure.
\item
  Both K9\_E model variants predict suppression ratio ≈ 2: additive
  (2.000 exactly), multiplicative (1.913). The predicted value ≈ 2 is
  not discriminating for current data quality.
\item
  P10-NOISE already established that 4 data points are insufficient to
  test any pattern --- A0B0 alone drives 80\% of Δχ².
\end{enumerate}

The path to pattern confirmation is K9-S12: an alpha-sweep protocol (QWP
angle 0°--90°) with N = 91,000 events and dedicated noise
characterization, providing direct 2BSM/1BSM comparison with sufficient
statistics.

\subsubsection{7.4 Adversarial Tests (4/4 PASS) and Operationalizability
Gates (3/3
PASS)}\label{adversarial-tests-44-pass-and-operationalizability-gates-33-pass}

Despite the empirically qualified status, K9\_E passes all structural
validity tests:

\textbf{Adversarial tests (4/4 PASS --- Phase 9):} 1. Born limit
recovery: β = 0 → P = Tr(E\_o ρ) exactly (verified analytically and
numerically). 2. Single-observer reduction: N = 1 → K\_ctx = ∅ → K9\_E =
Born (verified). 3. V-gate (Bhrānti): V = 0 events receive no
probability (verified). 4. isNull-gate (Anupalabdhi): null events (o =
∅) receive no probability (verified).

\textbf{Operationalizability gates (3/3 PASS --- G1/G2/G3, all 5.0/5):}
- G1: K9\_E is operationally distinguishable from Standard QM (δ\_S ≠ 0
when β \textgreater{} 0 and K\_ctx ≠ ∅): PASS. - G2: K9\_E has a
predeclared falsification condition (K9-S12 protocol, §8): PASS. - G3:
K9\_E's free parameter β is measurable in principle (Proietti-type EWF
experiment): PASS.

\textbf{Important boundary:} Adversarial tests and operationalizability
gates confirm that K9\_E is \emph{structurally consistent and testable}.
They do not confirm it is empirically correct. The 4/4 PASS result
establishes that K9\_E is a well-formed scientific hypothesis; the
P10-NOISE FAIL establishes that current empirical evidence is
insufficient to confirm it.

\begin{center}\rule{0.5\linewidth}{0.5pt}\end{center}

\subsection{8. K9-S12 Modified Bong
Protocol}\label{k9-s12-modified-bong-protocol}

This section summarizes the experimental specification for the first
dedicated test of the K9\_E probability postulate. The full protocol is
presented in the experimental child paper (paper\_002, arXiv submitted
2026-05-27). Section §8 provides the VVV-QMRF-level framing; the
photonic implementation details are in paper\_002.

\subsubsection{8.1 The Equatorial Gap: Why No Prior Test
Exists}\label{the-equatorial-gap-why-no-prior-test-exists}

\textbf{Proposition 1 (Equatorial Cancellation Theorem, paper\_002 §3):}
Every published optical EWF experiment --- Proietti et al.~(2019) and
Bong et al.~(2020) --- used Superobserver measurement settings at the
Bloch sphere equatorial plane (θ = π/2). At θ = π/2, the squared basis
overlaps \textbar⟨b\textbar d⟩\textbar² = 1/2 for all outcome pairs (b,
d). This makes the overlap-dependent suppression factor f\_perp constant
across all outcomes --- the suppression cancels in the normalization
Z\_E and P\_K9E reduces exactly to the Born rule. This is not a special
property of K9\_E; it holds for \textbf{any} overlap-only deformation of
the form P′ = P\_QM · g(\textbar⟨b\textbar d⟩\textbar²) / Z.

\textbf{Consequence:} Both published optical EWF experiments are
geometrically blind to the entire overlap-only class of deformations,
including K9\_E. No published experiment has varied the polar angle θ.
The K9-S12 protocol is the first proposal to break this fixed point.

\textbf{Signal structure:} The K9\_E signal δ⟨AB⟩(θ) vanishes
identically at θ = π/2 and is generically non-zero otherwise:

\begin{verbatim}
f_perp(+1, H) − f_perp(−1, H) = −cos θ                           [Eq. 11, paper_002]

At θ = π/2:  cos(π/2) = 0  →  f_perp constant  →  δ⟨AB⟩ = 0     [equatorial null]
At θ = 31°:  cos(31°) ≈ 0.857  →  f_perp outcome-dependent  →    [non-zero signal]
             δ⟨AB⟩ ≠ 0 for β > 0
\end{verbatim}

This vanishing-at-equator structure is a \textbf{genuine observable}
(Lemma 1, paper\_002 §3.2): it cannot be eliminated by any unitary
relabeling of the Superobserver's measurement basis.

\subsubsection{8.2 Single-QWP Design and
Predictions}\label{single-qwp-design-and-predictions}

\textbf{Experimental modification:} A single quarter-wave plate (QWP) is
re-inserted into the Bong et al.~(2020) apparatus, tilting the
Superobserver measurement from θ = π/2 to θ = 31°. No other hardware
change is required. N = 91,000 coincidences per setting (same as Bong
2020).

\textbf{Why θ = 31° (near-optimal):} Grid search over θ ∈ {[}0°, 90°{]}
shows θ = 35° yields FOM = 8.8 vs.~8.6 at θ = 31°. The broad plateau
means the exact optimum is non-critical: θ = 31° is near-optimal, not
uniquely optimal (paper\_002 v94 update). The key criterion is that θ ≠
π/2 --- any angle away from the equator probes the signal.

\textbf{Quantitative predictions at θ = 31° (exact numerical
computation):}

{\def\LTcaptype{none} % do not increment counter
\begin{longtable}[]{@{}
  >{\raggedright\arraybackslash}p{(\linewidth - 6\tabcolsep) * \real{0.2500}}
  >{\raggedright\arraybackslash}p{(\linewidth - 6\tabcolsep) * \real{0.2500}}
  >{\raggedright\arraybackslash}p{(\linewidth - 6\tabcolsep) * \real{0.2500}}
  >{\raggedright\arraybackslash}p{(\linewidth - 6\tabcolsep) * \real{0.2500}}@{}}
\toprule\noalign{}
\begin{minipage}[b]{\linewidth}\raggedright
Observable
\end{minipage} & \begin{minipage}[b]{\linewidth}\raggedright
Prediction
\end{minipage} & \begin{minipage}[b]{\linewidth}\raggedright
Significance
\end{minipage} & \begin{minipage}[b]{\linewidth}\raggedright
Role
\end{minipage} \\
\midrule\noalign{}
\endhead
\bottomrule\noalign{}
\endlastfoot
Gen LF 1 (Genuine LF Facet 1) & \textbf{+0.0891} & \textbf{8.6σ} &
Confirms LF violation is preserved at θ = 31° \\
δ⟨A₁B₂⟩ = ⟨A₁B₂⟩\emph{\{θ=31°\} − ⟨A₁B₂⟩}\{θ=π/2\} & \textbf{−0.0355} &
\textbf{20.8σ} & Core K9\_E signal --- non-zero iff K9\_E with β
\textgreater{} 0 \\
Minimum detectable β at 5σ & \textbf{0.07} & --- & Search sensitivity \\
Figure of Merit (FOM) & \textbf{8.6} & --- & Combined LF + K9\_E
signal \\
\end{longtable}
}

\textbf{Scope boundary (paper\_002 §2.3):} These predictions assume the
overlap-only class (Level 0). The experiment tests whether nature
exhibits overlap-dependent deformation; it does not test the full K9\_E
framework. Broader deformation classes (Level 1:
density-matrix-dependent; Level 2: multi-partite; Level 3:
non-geometric) lie outside Proposition 1's scope and require independent
experimental designs.

\textbf{Loophole status:} Under fair-sampling (η ≈ 0.87), K9-S12 is a
\textbf{loophole-open screening test}. A positive result (δ⟨A₁B₂⟩ ≠ 0 at
θ = 31°) requires independent verification including θ-sweeps
(paper\_002 §8.2). A null δ⟨AB⟩ across a full θ-sweep would falsify the
overlap-only class.

\textbf{Falsification conditions:} K9\_E (overlap-only class) is
falsified if δ⟨AB⟩ = 0 at θ = 31° within predeclared statistical
tolerance, confirmed by θ-sweep.

\subsubsection{8.3 Why IBM Quantum Was
Rejected}\label{why-ibm-quantum-was-rejected}

The IBM Quantum approach to K9\_E testing was evaluated and rejected
(v31 RCA 4.92/5) on the basis of a \textbf{double category error}:

\textbf{Error 1 (registration structure):} K9\_E requires a K-space
registration structure --- tuples k = ⟨M, o, cert, t, V⟩ with K3
self-certification (cert = 1 intrinsic), K4 default validity (V = 1 upon
admission), K5 invalidation (V → 0 via cross-registration
contradiction), and K7 closure (V\_prov → V\_final). Gate-model QPUs
implement unitary circuits with binary qubit states. There is no concept
of a K-side registration tuple, K\_R carrier set, or validity field V ∈
\{0,1\} on a QPU. K-space does not exist as a computational object on
IBM Quantum hardware.

\textbf{Error 2 (EWF scenario):} K9\_E specifically requires an Extended
Wigner's Friend configuration --- a Friend observer registering a
definite outcome within a sealed laboratory, and a Superobserver
performing an interference measurement on the entire Friend+system
composite. Gate-model circuits simulate quantum evolution but do not
implement the structural architecture of multiple registering systems
with K\_joint construction and cross-registration authority. The
``Friend'' and ``Superobserver'' roles in K9\_E are K-space structural
roles, not qubit register labels.

\textbf{Conclusion:} K9\_E's testable predictions are predictions about
photonic EWF statistics. Only an optical EWF apparatus implementing the
K9-S12 protocol (or equivalent) constitutes a valid test platform.

\subsubsection{8.4 Relationship to
paper\_002}\label{relationship-to-paper_002}

paper\_002 (``Have Optical Wigner's Friend Experiments Been Blind to a
Geometric Degree of Freedom?'') was submitted to arXiv 2026-05-27 as an
independent experimental child paper. The VVV-QMRF framing in §8.1--8.3
provides the K-space motivation for the experimental protocol; the
photonic implementation details, statistical analysis, and Supplemental
materials are in paper\_002.

paper\_002 is self-contained as a physics paper: it presents the
Equatorial Cancellation Theorem (Proposition 1) and single-waveplate
protocol without requiring the VVV-QMRF K-space framework. The K9\_E
postulate (§5) provides one theoretical motivation for the
overlap-dependent deformation parameter β; paper\_002 frames β as a
phenomenological search parameter analogous to Standard Model Extension
(SME) coefficients.

\begin{center}\rule{0.5\linewidth}{0.5pt}\end{center}

\subsection{9. Connections to Baumann-Brukner and
Frauchiger-Renner}\label{connections-to-baumann-brukner-and-frauchiger-renner}

This section documents how two specific papers in the EWF literature
connect to VVV-QMRF's K-space structure. These are \textbf{Layer 4
(Class D) structural compatibility assessments} --- they show that
VVV-QMRF's machinery can be applied to the scenarios in these papers,
not that the papers empirically confirm K9\_E.

\subsubsection{9.1 Baumann \& Brukner (2024): T\_BB Class C
Conditional}\label{baumann-brukner-2024-t_bb-class-c-conditional}

\textbf{Paper reference:} Baumann, V. and Brukner, Č. (2024). ``Wigner's
friend as a rational agent.'' {[}arXiv reference; fit plan source:
\texttt{09\_Fitting\_Baumann\_Brukner/BB\_VVV\_fit\_plan.md} v1.4{]}

\textbf{Scenario:} A Friend F measures a quantum system inside a sealed
lab; a Wigner W performs either a coherent (interference) measurement or
a projective (read-out) measurement on the lab. B\&B show that a
rational agent (W) cannot consistently assign memory to F after a
coherent measurement --- the ``no-awareness'' result.

\textbf{VVV-QMRF compatibility analysis
(BB\_VVV\_compatibility\_section.md v2.1):}

{\def\LTcaptype{none} % do not increment counter
\begin{longtable}[]{@{}
  >{\raggedright\arraybackslash}p{(\linewidth - 4\tabcolsep) * \real{0.3333}}
  >{\raggedright\arraybackslash}p{(\linewidth - 4\tabcolsep) * \real{0.3333}}
  >{\raggedright\arraybackslash}p{(\linewidth - 4\tabcolsep) * \real{0.3333}}@{}}
\toprule\noalign{}
\begin{minipage}[b]{\linewidth}\raggedright
Verification
\end{minipage} & \begin{minipage}[b]{\linewidth}\raggedright
Status
\end{minipage} & \begin{minipage}[b]{\linewidth}\raggedright
Finding
\end{minipage} \\
\midrule\noalign{}
\endhead
\bottomrule\noalign{}
\endlastfoot
V1: K5 ↔ B\&B q₀₀ \textless{} 0 condition & ⚠️ \textbf{PARTIAL (F4
triggered)} & R\_BB (near-readout, x ≈ 0) ≠ R\_K5 (interference, x ≈
π/4). B\&B's no-valid-joint-model fires near readout; K5 fires at
interference. Structurally \textbf{different failure modes} --- not
mathematically equivalent. F4 finding: at x = π/4 (maximum
interference), q₀₀ = 0.5 \textgreater{} 0 always, while
requires\_K\_joint = 1 and K5 fires. \\
V2: K7 closure ↔ B\&B memory change Δp & ✅ \textbf{PASS} & K7 closure
magnitude Δp = \textbar1 − 2α²\textbar{} · sin²(2x) / 2 matches B\&B
memory change with identical functional form. Verified at 5 test points
including readout (x≈0), maximal interference (x=π/4), and transition
region. \\
T\_BB (Option A) & ✅ \textbf{Class C (conditional)} & No-awareness
result derived via K5 → K6 → V=0 chain using K7\_trace + D\_enc. G1
CLOSED by K7\_trace §18 + D\_enc §19. Computational verification PASS
(\texttt{bb\_vvv\_t\_bb\_verification.py}). \\
T\_BB' (Option C) & ✅ \textbf{CLOSED (superseded)} & V1 F4 finding
falsifies T\_BB' Step 1 premise (R\_BB = R\_K5 was required). T\_BB
(Option A) computationally verified; T\_BB' superseded. \\
\end{longtable}
}

\textbf{P2-C (π/8 first-principles):} The compatibility check at angle x
= π/8 yields Δp = \textbar1 − 2α²\textbar{} / 4 exactly --- a
first-principles derivation from K7 closure (no free parameters). This
is \textbf{exact} for all α² ≠ 0.5 (symmetric state degeneracy at α² =
0.5 confirmed as expected: Δp = 0).

\textbf{T\_BB derivation chain:}

\begin{verbatim}
K5_prospective: W's coherent measurement → requires_K_joint(K_F, K_W) = 1
K7 closure:     K_F closes at t_close(K_F)
K7_trace:       Δ_closure(k_F) := V_prov(k_F) − V_final(k_F) ∈ {0,1}
D_enc:          Enc(M_aware, k_F) = 1 iff o(M_aware) counterfactually depends on Δ_closure
K5 fires:       requires_K_joint(M_aware, M_W) = 1 → V(M_aware) → 0
Result:         Friend F cannot retain awareness (V=0) → T_BB no-awareness result
\end{verbatim}

\textbf{Key boundary:} V1's F4 finding does not invalidate T\_BB. T\_BB
uses K7\_trace + D\_enc (Layer 2), not the R\_BB = R\_K5 equivalence.
What F4 establishes is that B\&B's no-valid-joint-model condition and
VVV-QMRF's K5 incommensurability capture \textbf{structurally different
aspects} of the Wigner's Friend problem --- they are compatible
frameworks analyzing different failure modes.

\subsubsection{9.2 Frauchiger \& Renner (2018): FR Paradox Avoided via
K5
V\_prov}\label{frauchiger-renner-2018-fr-paradox-avoided-via-k5-v_prov}

\textbf{Paper reference:} Frauchiger, D. and Renner, R. (2018).
``Quantum theory cannot consistently describe the use of itself.''
\emph{Nature Communications} 9, 3711.

\textbf{Scenario:} Four agents (F1, F2, W1, W2) in nested EWF setups. FR
show that under three assumptions (Q: universality, C: consistency, S:
single-outcome), agents reach mutually contradictory predictions. The
contradiction requires that V(k\_F1) = 1 AND V(k\_W1) = 1 simultaneously
--- both the Friend's registration and the Wigner's registration are
valid at the same time under the same cross-observer validity demand.

\textbf{VVV-QMRF FR avoidance chain:}

\begin{verbatim}
FR contradiction requires:   V(k_F1) = 1  AND  V(k_W1) = 1  simultaneously
                                              ↓
K5 fires at t₂:              K_F1 ⊥_K K_W1 within comparison context C_K
                              (W1's coherent measurement ≡ requires_K_joint = 1)
                                              ↓
K6 Authority:                 Auth(K_W1 → K_F1) = 1
                              → V(k_F1) → 0  (K5 post-closure irreversible)
                                              ↓
Premise fails:                V(k_F1) = 1 AND V(k_W1) = 1 NEVER simultaneously satisfied
                                              ↓
FR contradiction:             AVOIDED — K5 gates the joint validity condition at source
\end{verbatim}

\textbf{VVV-QMRF recast of FR assumption (S):} The FR argument requires
F1's outcome and W1's outcome to belong to the same validity domain
(joint cross-observer fact). K5 ⊥\_K prevents this: they are
structurally incommensurable registrations. VVV-QMRF does not
\emph{reject} (S); it \emph{scopes} (S) to within a single K-context.
The FR contradiction never assembles within VVV-QMRF's structural
machinery because the joint validity demand is blocked by K5 before it
can be evaluated.

\textbf{V\_FR2 verification (script-verified, 2026-05-28):} K7\_trace
records Δ\_closure(k\_F1) := V\_prov(k\_F1) − V\_final(k\_F1) at W1's
coherent measurement. This is \textbf{structurally identical} to T\_BB
Step 2 in B\&B --- the same K7\_trace definition applies unchanged to a
different scenario with different agents and a different contradiction
mechanism. \texttt{fr\_vvv\_k7trace\_consumer\_verification.py} v1.0:
OVERALL PASS.

\textbf{T\_FR status (Class D, blocked by G\_FR2):} A full T\_FR
(No-Joint-Validity Bridge Theorem) for the 4-agent scenario (F1, F2, W1,
W2) requires the N=4 K\_joint construction, which depends on T4-H for
N=4. T4-H is THEOREM for N=2 (T1 constructive) and general N (abstract
colimit proof, Steps 1-4). The N=4 concrete instantiation has not been
explicitly verified. A simplified 2-agent version (F1+W1) is feasible
and structurally established by V\_FR2.

\subsubsection{9.3 K7\_trace + D\_enc: Canonical Layer 2 with Four
Consumers}\label{k7_trace-d_enc-canonical-layer-2-with-four-consumers}

The V\_FR2 result confirms K7\_trace as a \textbf{scenario-agnostic}
Layer 2 primitive --- it applies with identical definition across
structurally distinct scenarios:

{\def\LTcaptype{none} % do not increment counter
\begin{longtable}[]{@{}
  >{\raggedright\arraybackslash}p{(\linewidth - 6\tabcolsep) * \real{0.2500}}
  >{\raggedright\arraybackslash}p{(\linewidth - 6\tabcolsep) * \real{0.2500}}
  >{\raggedright\arraybackslash}p{(\linewidth - 6\tabcolsep) * \real{0.2500}}
  >{\raggedright\arraybackslash}p{(\linewidth - 6\tabcolsep) * \real{0.2500}}@{}}
\toprule\noalign{}
\begin{minipage}[b]{\linewidth}\raggedright
Consumer
\end{minipage} & \begin{minipage}[b]{\linewidth}\raggedright
Scenario
\end{minipage} & \begin{minipage}[b]{\linewidth}\raggedright
Role
\end{minipage} & \begin{minipage}[b]{\linewidth}\raggedright
Status
\end{minipage} \\
\midrule\noalign{}
\endhead
\bottomrule\noalign{}
\endlastfoot
\textbf{T\_BB} (BB-VVV) & Baumann-Brukner (2024): F + W angle sweep &
V\_prov substitute after K7 closure for T\_BB no-awareness derivation &
Class C conditional \\
\textbf{D\_enc} & All scenarios with K7 closure & Parent: Enc(M\_aware,
k\_F) = 1 iff counterfactual dependence on Δ\_closure & Class
C-canonical \\
\textbf{3-OBS} & Three-observer hierarchical EWF & Δ\_closure
propagation through F1→F2→W hierarchy (§10) & Class C (T4-H verified) \\
\textbf{FR} (FR-VVV) & Frauchiger-Renner (2018): 4-agent nested EWF &
V\_FR2: Δ\_closure(k\_F1) after W1's coherent measurement --- blocks FR
joint validity & Class D (T\_FR blocked by G\_FR2) \\
\end{longtable}
}

\textbf{Architectural significance:} K7\_trace being consumed by four
structurally distinct analyses (B\&B angle-sweep, D\_enc semantic, 3-OBS
hierarchical, FR nested) confirms it is a legitimate Layer 2 primitive
--- not a single-use BB-VVV construct. This cross-validation was the
primary criterion for canonical promotion (RCA 4.77/5, 2026-05-27).

\begin{center}\rule{0.5\linewidth}{0.5pt}\end{center}

\subsection{10. T4-H THEOREM and 3-Observer
Prediction}\label{t4-h-theorem-and-3-observer-prediction}

\subsubsection{10.1 T4-H: From Hypothesis to Theorem (4/4 Steps,
2026-05-28)}\label{t4-h-from-hypothesis-to-theorem-44-steps-2026-05-28}

T4-H is the N-observer Colimit Existence Theorem for the category
C\_\{K-space\}. It states: \emph{for any finite diagram D of K-spaces
with K1--K8-preserving morphisms, the colimit K\_colim exists and
satisfies K1--K8.}

T4-H was originally framed as a hypothesis (``T4-H'') because its proof
required four sequential steps. As of 2026-05-28 all four are verified
(RCA 4.74/5):

{\def\LTcaptype{none} % do not increment counter
\begin{longtable}[]{@{}
  >{\raggedright\arraybackslash}p{(\linewidth - 6\tabcolsep) * \real{0.2500}}
  >{\raggedright\arraybackslash}p{(\linewidth - 6\tabcolsep) * \real{0.2500}}
  >{\raggedright\arraybackslash}p{(\linewidth - 6\tabcolsep) * \real{0.2500}}
  >{\raggedright\arraybackslash}p{(\linewidth - 6\tabcolsep) * \real{0.2500}}@{}}
\toprule\noalign{}
\begin{minipage}[b]{\linewidth}\raggedright
Step
\end{minipage} & \begin{minipage}[b]{\linewidth}\raggedright
Claim
\end{minipage} & \begin{minipage}[b]{\linewidth}\raggedright
Key technique
\end{minipage} & \begin{minipage}[b]{\linewidth}\raggedright
Status
\end{minipage} \\
\midrule\noalign{}
\endhead
\bottomrule\noalign{}
\endlastfoot
Step 1 & C\_\{K-space\} is a valid category & Identity morphisms +
composition + associativity via K8-preserving maps &
\textbf{VERIFIED} \\
Step 2 & K\_colim = (∐\_i K\_i)/\textasciitilde{} is constructible &
Lexicographic t-assignment (SP1), V-snapshot at embedding (SP2),
transitive closure for \textless\_colim (SP3); 5/5 gates PASS &
\textbf{VERIFIED} (RCA 4.73/5) \\
Step 3 & K\_colim satisfies K1--K8 & \textbf{T-PRES Lemma} (all
morphisms preserve t exactly, resolving K2 acyclicity); K3--K8 verified
via quotient structure & \textbf{VERIFIED} (RCA 4.74/5, 2026-05-28) \\
Step 4 & K\_colim satisfies universal property & u({[}k,i{]}) := f\_i(k)
well-defined (T-REP corollary), K8-preserving, unique &
\textbf{VERIFIED} (2026-05-28) \\
\end{longtable}
}

\textbf{T-PRES Lemma (key to Step 3):} Any morphism f: K\_i → K\_j in
C\_\{K-space\} satisfies t\_j(f(k)) = t\_i(k) for all k ∈ K\_i. This is
because K8-preservation requires all five tuple fields (M, o, cert, t,
V) to be preserved exactly. Consequence: all representatives of any
equivalence class {[}k, i{]} share the same t value (T-REP Corollary),
ensuring K2's acyclicity condition holds in K\_colim without cycles.

\textbf{Claim class:} T4-H is now a \textbf{full THEOREM}, not a
hypothesis or conditional claim. T1 (N=2 constructive) remains
independently valid without invoking T4-H's universal property. T4, T5
(K\_joint Associativity) are upgraded accordingly.

\subsubsection{10.2 Upgrade: 3-Observer Prediction from Class
C-Conditional to Class
C}\label{upgrade-3-observer-prediction-from-class-c-conditional-to-class-c}

The 3-observer registration transition mechanism
(\texttt{3observer\_registration\_transition.md}, v1.1, 2026-05-28) was
previously classified \textbf{Class C-conditional} pending T4-H Steps
3-4. With T4-H THEOREM confirmed, the conditional gate is resolved.

\textbf{3-observer hierarchical configuration:}

\begin{verbatim}
Hierarchy: t_F1 < t_F2 < t_W

  F1 measures system S → registers k_F1 in K_{F1}
  F2 measures F1's lab → registers k_F2 in K_{F2}
  W (Superobserver) measures joint F1+F2 lab
     (requires_K_joint = 1; interference measurement)
\end{verbatim}

The registration transition mechanism traces via four steps:

\begin{itemize}
\tightlist
\item
  \textbf{H1:} K7\_trace records Δ\_closure(k\_F1) = V\_prov(k\_F1) −
  V\_final(k\_F1) at t\_close(K\_\{F1\}).
\item
  \textbf{H2:} K7\_trace records Δ\_closure(k\_F2) similarly at
  t\_close(K\_\{F2\}).
\item
  \textbf{H3:} D\_enc evaluates: Enc(M\_aware, k\_F1) = 1 iff the
  post-closure act M\_aware's outcome depends counterfactually on
  Δ\_closure(k\_F1). In the 3-OBS case, this propagation requires
  K\_joint for N=3, provided by T4-H.
\item
  \textbf{H4:} If Enc = 1 → requires\_K\_joint(M\_aware, M\_W) = 1 → K5
  fires → V(M\_aware) → 0 → no-awareness result propagates
  hierarchically.
\end{itemize}

\textbf{Upgraded classification:}

{\def\LTcaptype{none} % do not increment counter
\begin{longtable}[]{@{}
  >{\raggedright\arraybackslash}p{(\linewidth - 6\tabcolsep) * \real{0.2500}}
  >{\raggedright\arraybackslash}p{(\linewidth - 6\tabcolsep) * \real{0.2500}}
  >{\raggedright\arraybackslash}p{(\linewidth - 6\tabcolsep) * \real{0.2500}}
  >{\raggedright\arraybackslash}p{(\linewidth - 6\tabcolsep) * \real{0.2500}}@{}}
\toprule\noalign{}
\begin{minipage}[b]{\linewidth}\raggedright
Component
\end{minipage} & \begin{minipage}[b]{\linewidth}\raggedright
Prior class
\end{minipage} & \begin{minipage}[b]{\linewidth}\raggedright
Current class
\end{minipage} & \begin{minipage}[b]{\linewidth}\raggedright
Condition
\end{minipage} \\
\midrule\noalign{}
\endhead
\bottomrule\noalign{}
\endlastfoot
K7\_trace + D\_enc application & C-conditional & \textbf{Class C} & T4-H
THEOREM 2026-05-28 \\
No-awareness propagation (H1-H4) & C-conditional & \textbf{Class C} &
T4-H THEOREM 2026-05-28 \\
Numerical prediction δ\_M3 & Illustrative & Illustrative & Conditional
on K9\_E P9 only (see §10.3) \\
\end{longtable}
}

\subsubsection{10.3 3-Observer Numerical Prediction
(Illustrative)}\label{observer-numerical-prediction-illustrative}

\textbf{Important scope boundary:} The 3-observer \emph{structural
mechanism} (§10.2) is Class C --- it follows from K-space axioms K1--K8
+ Layer 2 (K7\_trace, D\_enc) + T4-H THEOREM. The 3-observer
\emph{numerical prediction} (δ\_M3) requires the additional K9\_E
postulate (P9) --- it is illustrative and conditional on K9\_E.

\textbf{Prediction (conditional on K9\_E P9):}

\begin{verbatim}
δ_M3 = −0.223  at β = 0.3  (illustrative, not a measurement)
\end{verbatim}

\begin{itemize}
\tightlist
\item
  3-observer K\_ctx is larger than 2-observer K\_ctx → amplification
  factor ≈ 11× relative to 2-observer δ\_S.
\item
  At β = 0.3, the 2-observer prediction gives δ\_S ≈ −0.020; the
  3-observer setup predicts δ\_M3 ≈ −0.223.
\item
  The amplification arises because K\_ctx(k\_i, Exp) includes two
  contextual observer registrations instead of one, increasing f\_perp.
\end{itemize}

\textbf{Dependency chain for numerical prediction:}

\begin{verbatim}
K7_trace + D_enc (Layer 2, canonical)
    ↓ mechanism (Class C)
T4-H THEOREM (N=3 K_joint — resolved)
    ↓ structural substrate
K9_E postulate P9 (Class C qualified — not yet confirmed)
    ↓ probability assignment
δ_M3 = −0.223 at β=0.3 (illustrative)
\end{verbatim}

\textbf{Caveat:} β = 0.3 is chosen for illustration (not from a fit to
3-observer data, which does not yet exist). The true β value, if K9\_E
is correct, would be determined by a future 3-observer experiment. The
11× amplification factor itself is structural (from K\_ctx size) and
holds for any non-zero β, not just β = 0.3.

\begin{center}\rule{0.5\linewidth}{0.5pt}\end{center}

\subsection{11. φ Conjecture: Structure-Preserving Map K → B(H) (Class D
Supporting)}\label{ux3c6-conjecture-structure-preserving-map-k-bh-class-d-supporting}

\begin{quote}
\textbf{This section is Class D supporting material --- NOT the central
claim of v3.0.} The central claims of this paper are the K-space
architecture (§4), the K9\_E postulate (§5), and the K9-S12 experimental
specification (§8). The φ conjecture (Project B) is logically
independent of K9\_E (Project C) --- see §5.6 for the independence
statement.
\end{quote}

\subsubsection{11.1 The Conjecture and Its
Motivation}\label{the-conjecture-and-its-motivation}

\textbf{Conjecture (Class D):} There exists a structure-preserving map
φ: K → B(H), where B(H) is the algebra of bounded operators on Hilbert
space H, such that φ encodes the registration-logic properties of
K-space (K1--K8) in the operator-algebraic language of Standard QM.

\textbf{Motivation:} VVV-QMRF's K-space axiomatizes \emph{what gets
registered} (registration-logic structure). Standard QM's B(H)
axiomatizes \emph{what can be measured} (operator-algebraic structure).
If φ exists, it would formalize the correspondence between registration
events in K and observable operators in B(H) --- not by identifying K
with H, but by showing that K's registration-logic has a faithful
structural image in operator-algebraic language.

\textbf{Phase 1--4 completion (Project B Track B):} The φ-map research
program has completed four phases: - \textbf{Phase 1:} Target selection
--- B(H) chosen as codomain with image Im(φ) ⊆ \{P\_o\} ∪ \{0\}
(projection sub-lattice). Three counter-arguments against B(H) resolved.
- \textbf{Phase 2:} Necessary conditions N₁--N\_T derived from K1--K8
and T1--T9 (§11.2). - \textbf{Phase 3:} φ-conditional interpretation
scope analysis (§12.2). - \textbf{Phase 4:} Open items documented (φ-O2:
N₆ sufficiency boundary confirmed as fundamental; C2 readiness 8.0/10).

\textbf{Why B(H) as codomain (not C*-algebra, von Neumann algebra, or
projection lattice):}

{\def\LTcaptype{none} % do not increment counter
\begin{longtable}[]{@{}
  >{\raggedright\arraybackslash}p{(\linewidth - 2\tabcolsep) * \real{0.5000}}
  >{\raggedright\arraybackslash}p{(\linewidth - 2\tabcolsep) * \real{0.5000}}@{}}
\toprule\noalign{}
\begin{minipage}[b]{\linewidth}\raggedright
Alternative
\end{minipage} & \begin{minipage}[b]{\linewidth}\raggedright
Problem
\end{minipage} \\
\midrule\noalign{}
\endhead
\bottomrule\noalign{}
\endlastfoot
P(H) projection lattice & Not an algebra --- K6 authority-composition
(φ-6) has no natural image \\
Von Neumann algebra M & Requires specifying which subalgebra --- adds
unwarranted structural choice \\
Functor to category C\_obs & Loses concrete B(H) operator intuition;
abstract for current phase \\
\textbf{B(H) (chosen)} & Contains all projections P\_o; admits φ-1
through φ-7' conditions; standard QM home \\
\end{longtable}
}

\textbf{EWF connection:} In the EWF scenario, K\_F ⊥\_K K\_W on the
K-side maps to {[}ι(P\_\{o\_F\}), P\_\{o\_W\}{]} ≠ 0 in B(H) --- the
familiar non-commutativity of quantum observables. The φ-conditional
analysis (§12.2) re-frames each interpretation's structural gap as the
specific necessary condition N\_i it lacks machinery to satisfy.

\subsubsection{11.2 Necessary Conditions N₁--N\_T (Derived from
K1--K8)}\label{necessary-conditions-nux2081n_t-derived-from-k1k8}

The following necessary conditions must hold for any map φ: K → B(H) to
qualify as structure-preserving. These are necessary conditions --- not
sufficient for φ's existence.

{\def\LTcaptype{none} % do not increment counter
\begin{longtable}[]{@{}
  >{\raggedright\arraybackslash}p{(\linewidth - 6\tabcolsep) * \real{0.2500}}
  >{\raggedright\arraybackslash}p{(\linewidth - 6\tabcolsep) * \real{0.2500}}
  >{\raggedright\arraybackslash}p{(\linewidth - 6\tabcolsep) * \real{0.2500}}
  >{\raggedright\arraybackslash}p{(\linewidth - 6\tabcolsep) * \real{0.2500}}@{}}
\toprule\noalign{}
\begin{minipage}[b]{\linewidth}\raggedright
Condition
\end{minipage} & \begin{minipage}[b]{\linewidth}\raggedright
Source axiom
\end{minipage} & \begin{minipage}[b]{\linewidth}\raggedright
Statement
\end{minipage} & \begin{minipage}[b]{\linewidth}\raggedright
B(H) encoding
\end{minipage} \\
\midrule\noalign{}
\endhead
\bottomrule\noalign{}
\endlastfoot
\textbf{φ-1} (Well-Definedness) & K1 & φ: K\_R → B(H) is a total
function; ∀k ∈ K\_R, φ(k) ∈ B(H) defined & cert(k) = 1 for all k ∈ K\_R
→ no undefined images \\
\textbf{φ-2} (Lüders Order) & K2 & k₁ \textless{}\emph{R k₂ maps to
Lüders update: ρ ↦ P}\{o₂\} · P\_\{o₁\} · ρ · P\_\{o₁\} · P\_\{o₂\}
(application order from temporal order) & Non-commutativity of
P\_\{o₁\}, P\_\{o₂\} mirrors non-symmetry of \textless\_R \\
\textbf{φ-3} (Cert-Reflection) & K3 & φ(k) is determined by k's own
tuple fields alone; for all R' ≠ R: φ\_R(k) ∉ φ\_\{R'\}(k') & P\_o =
\textbar o⟩⟨o\textbar{} determined entirely by the registered outcome o
--- intrinsic \\
\textbf{φ-4} (Validity-Positivity) & K4 & V(k) = 1 → φ(k) = P\_o ≥ 0,
φ(k) ≠ 0; V(k) = 0 → φ(k) = 0 & Connects to AOE (Bong 2020): V(k) = 1
(default validity) → non-zero positive operator \\
\textbf{φ-5} (Invalidation-Absorption) & K5 & V\_final(k) = 0
(post-closure irreversible) → φ(k) = 0; P\_o → 0 is irreversible in B(H)
& Zero operator is an absorbing element: 0·A = A·0 = 0; no recovery
without additional info \\
\textbf{φ-6} (Authority-Composition) & K6 & Auth(k₂ → k₁, C\_K) = 1 maps
to P\_\{o₂\} · P\_\{o₁\} ≠ 0 (non-annihilating composition) & k₂ has
authority over k₁ only if their projectors are not onto orthogonal
subspaces \\
\textbf{φ-7} (Embedding Naturality) & K8 & Diagram commutes: φ\_X(i(k))
= ι(φ\_R(k)) where ι: B(H\_R) → B(H\_X) is tensor extension & ι(P\_o) =
P\_o ⊗ 1\_\{H\_extra\} --- standard tensor product inclusion of
subsystem projector \\
\textbf{φ-7′} (Closure Finalization) & K7 & φ\_final(k) = φ\_prov(k) at
t\_close; post-closure φ is fixed; provisional φ\_prov can change
pre-closure & Fixes the operator image at t\_close --- after that, φ(k)
= 0 or P\_o cannot change \\
\textbf{N\_T} (Theorem bridge) & T2, T3 & K\_F ⊥\_K K\_W (K-side) →
{[}ι(P\_\{o\_F\}), P\_\{o\_W\}{]} ≠ 0 (B(H)-side): non-commutativity as
image of incommensurability & This is the φ-conditional EWF connection:
operator non-commutativity is the B(H) image of K-side ⊥\_K \\
\end{longtable}
}

\textbf{Derivation status:} N₁--N\_T are derived under the assumption
that φ exists. They state what φ \emph{must} satisfy; they do not prove
φ \emph{does} exist. The relationship is: φ exists → N₁--N\_T hold. The
converse (N₁--N\_T hold → φ exists) remains unproven. Open item φ-O2 (N₆
sufficiency for K6 authority-composition) is documented as a fundamental
boundary.

\subsubsection{11.3 K ≠ H Boundary
Reaffirmed}\label{k-h-boundary-reaffirmed}

A critical architectural boundary: \textbf{φ is a map from K to B(H),
not an identification K = H.}

Three distinct layers remain separate throughout the φ analysis: -
\textbf{K (registration layer):} Tuples ⟨M, o, cert, t, V⟩ with
registration-logic operations (K1--K8). Not a Hilbert space, not a
probability space. - \textbf{ρ (physical state layer):} Density
operators ρ ∈ D(H), evolving under Standard QM (P1--P4). Not modified by
VVV-QMRF. - \textbf{B(H) (observable algebra):} Bounded operators,
containing projection operators P\_o. The codomain of φ.

φ: K → B(H) maps registration events to operators, but it does not
collapse K into H. The image Im(φ) ⊆ \{P\_o\} ∪ \{0\} is a small subset
of B(H) --- φ selects the operators relevant to registration from within
the full operator algebra. This is the formal expression of K ≠ H:
K-side validity events map \emph{into} B(H) without being
\emph{identified with} H.

\textbf{Why this matters:} If K = H were asserted, VVV-QMRF would be
claiming that every registration event has a direct physical
correspondent in the Hilbert space --- this would constitute a
modification of Standard QM. K ≠ H (preserved by φ) ensures VVV-QMRF
remains a registration-layer analysis rather than a modification of
quantum dynamics.

\subsubsection{11.4 Class D Supporting
Status}\label{class-d-supporting-status}

The φ conjecture is Class D: - The \textbf{existence} of φ is not
proven. No constructive proof has been given that a map satisfying all
of φ-1 through φ-7', N\_T is well-defined for all K-spaces. - The
necessary conditions are \textbf{not sufficient}. Satisfying N₁--N\_T
would be required for φ to exist, but whether it is sufficient is an
open mathematical question. - \textbf{No empirical prediction} derives
solely from φ. K9\_E (§5) provides all testable predictions
independently of φ. - φ is \textbf{not falsified} by a K9\_E null
result. The φ conjecture concerns the formal relationship K → B(H), not
the probability assignment P(o\textbar K).

The value of the φ analysis is interpretive: it provides the formal
language for §12.2 (φ-conditional scope boundaries for existing QM
interpretations) and maps the VVV-QMRF framework into a notation
familiar to operator-algebraic quantum mechanics. It is supporting
infrastructure for the project's long-term formalization goal, not its
current testable claim.

\begin{center}\rule{0.5\linewidth}{0.5pt}\end{center}

\subsection{12. Positioning Against Existing
Interpretations}\label{positioning-against-existing-interpretations}

\subsubsection{12.1 Architectural
Comparison}\label{architectural-comparison}

{\def\LTcaptype{none} % do not increment counter
\begin{longtable}[]{@{}
  >{\raggedright\arraybackslash}p{(\linewidth - 4\tabcolsep) * \real{0.3333}}
  >{\raggedright\arraybackslash}p{(\linewidth - 4\tabcolsep) * \real{0.3333}}
  >{\raggedright\arraybackslash}p{(\linewidth - 4\tabcolsep) * \real{0.3333}}@{}}
\toprule\noalign{}
\begin{minipage}[b]{\linewidth}\raggedright
Interpretation
\end{minipage} & \begin{minipage}[b]{\linewidth}\raggedright
Response to EWF paradox
\end{minipage} & \begin{minipage}[b]{\linewidth}\raggedright
VVV-QMRF difference
\end{minipage} \\
\midrule\noalign{}
\endhead
\bottomrule\noalign{}
\endlastfoot
Copenhagen & WF is ill-posed; classical apparatus required & K3/K4
formalize why apparatus has registration authority: σ\_R(M) = 1
intrinsically (no Heisenberg cut needed) \\
Many-Worlds & All outcomes occur; no paradox & Registration events are
singular per K-side (K1 cert = 1 invariant); not globally branching \\
QBism & Facts are agent-relative; no paradox & Agrees on
agent-relativity; adds formal K-side structure (K1--K8, six conditions)
that QBism does not supply \\
Relational QM & Facts are relation-relative; no paradox & Closest
existing framework; VVV-QMRF adds formal machinery (K3 self-cert, K7
closure, requires\_K\_joint) that RQM does not \\
Objective Collapse (GRW) & Physical collapse restores absolute facts &
VVV-QMRF does not require physical collapse on ρ-side; K5 operates on
K-side validity, independent of ρ dynamics \\
\end{longtable}
}

\textbf{Relation to Relational Quantum Mechanics:} Rovelli (1996) argues
quantum states are relative to observers. VVV-QMRF adds the formal
registration-layer structure explaining \emph{why} they are relative:
σ\_R(M) operates intrinsically within K\_R, independently of K\_\{R′\}.
K\_F ⊥\_K K\_W is the registration-layer account of Rovelli's
observer-relativity. This is a genuine extension: VVV-QMRF supplies the
formal conditions (K3 self-cert, K2 temporal order, K5 invalidation, K7
closure, requires\_K\_joint predicate) that Relational QM does not.

\subsubsection{12.2 φ-Conditional Scope Boundaries (Class
D)}\label{ux3c6-conditional-scope-boundaries-class-d}

Section §11 derives necessary conditions N₁--N\_T for any
structure-preserving map φ: K → B(H). These conditions re-frame each
architectural gap in §12.1 as a \textbf{φ-conditional scope boundary}:
the specific condition N\_i that the interpretation lacks the structural
machinery to satisfy.

\textbf{Scope boundary convention:} ``Lacks the structural machinery for
N\_i'' means the interpretation does not supply the formal element
required by N\_i --- not that the interpretation is incorrect within its
own domain.

{\def\LTcaptype{none} % do not increment counter
\begin{longtable}[]{@{}
  >{\raggedright\arraybackslash}p{(\linewidth - 6\tabcolsep) * \real{0.2500}}
  >{\raggedright\arraybackslash}p{(\linewidth - 6\tabcolsep) * \real{0.2500}}
  >{\raggedright\arraybackslash}p{(\linewidth - 6\tabcolsep) * \real{0.2500}}
  >{\raggedright\arraybackslash}p{(\linewidth - 6\tabcolsep) * \real{0.2500}}@{}}
\toprule\noalign{}
\begin{minipage}[b]{\linewidth}\raggedright
Interpretation
\end{minipage} & \begin{minipage}[b]{\linewidth}\raggedright
Architectural gap (§12.1)
\end{minipage} & \begin{minipage}[b]{\linewidth}\raggedright
φ-conditional scope boundary
\end{minipage} & \begin{minipage}[b]{\linewidth}\raggedright
N\_i lacking
\end{minipage} \\
\midrule\noalign{}
\endhead
\bottomrule\noalign{}
\endlastfoot
Copenhagen & No formal definition of classical apparatus & Cannot
construct φ-domain element k = ⟨M, o, cert, t, V⟩ with cert field → φ
has no defined domain & \textbf{N₃} (cert-reflection, §11.2): cert =
σ\_R(M) = 1 required to admit k into K\_R \\
Many-Worlds & No physical observable distinguishes branches & No
singular admission act per branch → K\_R carrier set underspecified → φ
total function (N₁) has no well-defined domain & \textbf{N₁}
(well-definedness, §11.2): φ total requires K\_R well-defined; branching
leaves carrier set underspecified \\
QBism & Subjective probability is not a physical quantity & No
structural V(k) ∈ \{0,1\} field: agent-relative belief is not a binary
validity predicate → φ cannot satisfy validity-positivity & \textbf{N₄}
(validity-positivity, §11.2): V(k) = 1 → φ(k) = P\_o ≥ 0 requires V as
structural field, not degree of belief \\
Relational QM & No temporal closure event & RQM has relational facts but
no t\_close → φ\_final cannot be fixed & \textbf{N₇}
(closure-finalization, §11.2): φ = φ\_final fixed at t\_close requires
temporal closure boundary absent from RQM \\
Objective Collapse (GRW) & Collapse on ρ-side only & GRW collapse
operates on ρ-side state; K-side V: 1→0 irreversible transition not
formally supplied & \textbf{N₅} (invalidation-absorption, §11.2): K-side
V-update and post-closure irreversibility not entailed by ρ-side
collapse \\
\end{longtable}
}

\textbf{Status --- Class D:} N₁--N\_T are derived under the assumption
that φ: K → B(H) exists (§11). This §12.2 analysis means: \emph{if φ
were to exist, these are the structural elements existing
interpretations would need to supply.} Neither φ's existence nor the
empirical correctness of these scope boundaries is asserted. See §13.4
for the relationship between this φ-conditional analysis and the broader
architectural constraint framing.

\subsubsection{12.3 Three Logically Independent
Projects}\label{three-logically-independent-projects}

VVV-QMRF comprises three projects related by a one-way motivation chain.
This section makes that independence explicit to prevent the critique
that Buddhist epistemology is post-hoc justification for a physics
claim.

\textbf{Project A --- BE↔QM Comparative Mapping (interpretive
framework):} - 30 nodes (N\_BE\_00001--N\_BE\_00030), 39 edges
(ED\_BE\_00001--ED\_BE\_00039) - Comparative framework between Buddhist
Pramāṇa epistemology and quantum measurement concepts - Independently
testable as a comparative-philosophy framework (internal consistency,
historical accuracy of BE nodes, adequacy of QM node definitions) -
\textbf{Does NOT constitute evidence for K-space axioms or K9\_E}

\textbf{Project B --- VVV-QMRF Conceptual Framework (formal
axiomatization):} - K1--K8 frozen Layer 1 axioms; T1--T9 updatable Layer
2 bridge theorems; φ-map Class D conjecture - Independently testable as
a formal axiom system (consistency, non-redundancy, expressive power) -
Motivated by Project A's structural analysis of valid registration, but
not logically derived from it - \textbf{Does NOT make testable physical
predictions by itself}

\textbf{Project C --- K9\_E Testable Hypothesis (falsifiable claim):} -
K9\_E postulate P9 with one free parameter β; Class C (qualified) -
Falsifiable via K9-S12 Modified Bong Protocol (§8); predicts δ⟨A₁B₂⟩ ≠ 0
at θ = 31° - Motivated by Project B's K-space structure (K5\_prospective
→ T8 → f\_perp), but K9\_E is a \emph{postulate} not \emph{derivable}
from K1--K8 - \textbf{This is the only project with a current
falsification path}

\textbf{A→B→C: one-way motivation, not logical derivation:}

\begin{verbatim}
Project A (BE↔QM)  →(motivates)→  Project B (K-space)  →(motivates)→  Project C (K9_E)
    ↑                                    ↑                                    ↑
interpretive                          formal                             testable
framework                          axiomatization                       hypothesis
\end{verbatim}

Each project is independently falsifiable: - Project A is falsified if
the BE node definitions are historically inaccurate or the QM node
definitions mischaracterize Standard QM. - Project B is falsified if
K1--K8 are internally inconsistent, or if the K-space formalism cannot
be applied consistently to any EWF scenario. - Project C is falsified if
δ⟨AB⟩ = 0 at θ = 31° across a full θ-sweep (K9-S12 null result).

\textbf{A null K9\_E result (β = 0 measured by K9-S12) falsifies Project
C but does not invalidate Project B (the K-space architecture remains
valid as a formal structure) or Project A (the BE↔QM comparative
framework stands independently).}

\begin{center}\rule{0.5\linewidth}{0.5pt}\end{center}

\subsection{13. Scope, Limitations, and Open
Items}\label{scope-limitations-and-open-items}

\subsubsection{13.1 What This Paper Does Not
Claim}\label{what-this-paper-does-not-claim}

\begin{itemize}
\tightlist
\item
  Does not replace Standard Quantum Mechanics or revise P1--P4.
\item
  Does not revise the Born rule or unitary evolution (K9\_E recovers
  Born exactly at β = 0).
\item
  Does not claim Buddhist epistemology proves quantum mechanics (Project
  A is interpretive framework, not evidence).
\item
  Does not claim K-side incommensurability is experimentally confirmed.
\item
  Does not claim that consciousness plays any role in registration.
\item
  Does not claim the EWF paradox is fully resolved on the ρ-side.
\item
  Does not claim requires\_K\_joint is a necessary-and-sufficient
  condition for all LF violations; it is proposed only as a necessary
  registration-layer condition.
\item
  \textbf{Does not claim K9\_E is experimentally confirmed} --- K9\_E is
  Class C (qualified): structurally testable, empirically unconfirmed.
  The 2.31σ Proietti fit is real but ambiguous; noise cannot be ruled
  out as an alternative explanation (§7).
\item
  \textbf{Does not claim the φ-map φ: K → B(H) is proven to exist} ---
  the map is conjectured (Class D). The necessary conditions N₁--N\_T
  are derived but not sufficient for φ's existence (§11).
\item
  \textbf{Does not claim K9\_E is the unique probability rule consistent
  with K1--K8} --- K1--K8 are structural axioms that do not uniquely
  determine a probability rule; K9\_E is one motivated postulate, not
  the only possibility.
\item
  \textbf{Does not claim the 3-observer prediction δ\_M3 = -0.223 is
  independent of β} --- this value is illustrative at β = 0.3; the true
  β is unknown until experimental measurement.
\end{itemize}

\subsubsection{13.2 Formal Phase Summary and Open
Items}\label{formal-phase-summary-and-open-items}

The ⊥\_K formal definition chain is complete at the proposed Class D/C
level:

{\def\LTcaptype{none} % do not increment counter
\begin{longtable}[]{@{}llll@{}}
\toprule\noalign{}
Formal layer & Symbol & Class & Defined in \\
\midrule\noalign{}
\endhead
\bottomrule\noalign{}
\endlastfoot
K-side incommensurability & ⊥\_K & D & §6.4 \\
Admissible joint K-space & K\_joint / AdmJoint & D & §6.3 \\
Joint-validity demand & D\_joint / requires\_K\_joint & D & §6.3 \\
K-side comparison context & C\_K & D & §6.4 \\
Bridge lemma & Bridge\_EWF & D/C & §6.5 \\
Operational data criterion & ODC\_K & C & §6.6 \\
K-Space axioms (Layer 1) & K1--K8 & D (K1=C) & §4.1--4.2 \\
Bridge theorems (Layer 2) & T1--T9, K7\_trace, D\_enc & D/C &
§4.3--4.4 \\
Probability postulate & K9\_E (P9) & C (qualified) & §5 \\
Experimental protocol & K9-S12 & C & §8 \\
φ conjecture & φ: K → B(H) & D (supporting) & §11 \\
\end{longtable}
}

\textbf{Resolved items (updated from v2.0):}

{\def\LTcaptype{none} % do not increment counter
\begin{longtable}[]{@{}
  >{\raggedright\arraybackslash}p{(\linewidth - 4\tabcolsep) * \real{0.3333}}
  >{\raggedright\arraybackslash}p{(\linewidth - 4\tabcolsep) * \real{0.3333}}
  >{\raggedright\arraybackslash}p{(\linewidth - 4\tabcolsep) * \real{0.3333}}@{}}
\toprule\noalign{}
\begin{minipage}[b]{\linewidth}\raggedright
Item
\end{minipage} & \begin{minipage}[b]{\linewidth}\raggedright
v2.0 status
\end{minipage} & \begin{minipage}[b]{\linewidth}\raggedright
v3.0 status
\end{minipage} \\
\midrule\noalign{}
\endhead
\bottomrule\noalign{}
\endlastfoot
Axiomatize K as full mathematical structure & Deferred long-term task &
\textbf{RESOLVED} --- K1--K8 v2.4, T1--T9, K7\_trace, D\_enc (§4) \\
K5\_prospective (A1 assumption) & Class D semantic extension &
\textbf{RESOLVED} --- upgraded to K5\_prospective clause (§4.2) \\
T4-H colimit existence & Conditional hypothesis & \textbf{THEOREM} (4/4
steps, RCA 4.74/5, 2026-05-28, §10.1) \\
3-OBS Class C-conditional & Conditional on T4-H & \textbf{Class C} (T4-H
resolved, §10.2) \\
K7\_trace canonical status & BB-VVV local construct & \textbf{Canonical
Layer 2} (4 consumers, RCA 4.77/5, §4.4) \\
\end{longtable}
}

\textbf{Deferred proof items (not required for current claim class):}

{\def\LTcaptype{none} % do not increment counter
\begin{longtable}[]{@{}
  >{\raggedright\arraybackslash}p{(\linewidth - 2\tabcolsep) * \real{0.5000}}
  >{\raggedright\arraybackslash}p{(\linewidth - 2\tabcolsep) * \real{0.5000}}@{}}
\toprule\noalign{}
\begin{minipage}[b]{\linewidth}\raggedright
Deferred item
\end{minipage} & \begin{minipage}[b]{\linewidth}\raggedright
Reason deferred
\end{minipage} \\
\midrule\noalign{}
\endhead
\bottomrule\noalign{}
\endlastfoot
Full semantic proof for Bridge\_EWF & Requires formal semantics of ``no
admissible reinterpretation''; paper uses operational sufficient
conditions \\
Full formal proof for ⊥\_K as mathematical relation & K1--K8 structural
definition complete; full topological/order-theoretic treatment
deferred \\
AdmJoint necessary-and-sufficient conditions & Currently sufficient
conditions A--E; full characterization is an open item \\
Equivalence of σ(M) and R̂\_svasa formalisms & Separate research track \\
\end{longtable}
}

\subsubsection{13.3 Falsification Condition
(Restated)}\label{falsification-condition-restated}

\textbf{K9\_E (Project C) is falsified} by: a K9-S12 experimental result
with δ⟨A₁B₂⟩ = 0 at θ = 31° within predeclared statistical tolerance
across a full θ-sweep. This would falsify the overlap-only class of
deformations (the class that K9\_E belongs to), not just a specific β
value.

\textbf{The K-side incommensurability conjecture is falsified} by: an
EWF experiment producing results consistent with a single K\_joint
registration model satisfying Conditions 1--6 for both observers
simultaneously, for a configuration independently classified as
requires\_K\_joint = 1 via D\_joint and Bridge\_EWF. If K\_joint exists
empirically, the Bridge\_EWF conditional bridge requires revision.

\textbf{K9\_E falsification (K9-S12 null result) does not falsify:} -
The K-space architectural framework (K1--K8 + Layer 2) --- Project B -
The K-side incommensurability conjecture (K\_F ⊥\_K K\_W) --- a
different structural claim - The BE↔QM comparative mapping --- Project A

Operationally, all three falsification tests require a predefined data
criterion (tolerance τ fixed before data collection); otherwise the
criterion becomes post-hoc.

\subsubsection{13.4 Architectural Constraints of the K-side/ρ-side
Separation}\label{architectural-constraints-of-the-k-sideux3c1-side-separation}

The core architectural commitment of VVV-QMRF is K ≠ H: the registration
layer is structurally distinct from the physical layer (§2.1). This
separation enables the framework to address the registration layer gap.
It also imposes two architectural constraints that are inherent
consequences of the layer separation, not oversights.

\textbf{Constraint 1 --- Self-certification is not a ρ-side observable.}
Condition 4 requires σ\_R(M) = 1 to be determined intrinsically within
K\_R. Any physical measurement of σ would itself require a second-order
registration act, re-entering the von Neumann regress that K3 is
designed to terminate. The operational compromise is stated in ODC\_K
condition (ii): self-certification is treated as a structural postulate
satisfied by construction when the registering system meets the
architectural definition of a K-side observer.

\textbf{Constraint 2 --- D\_joint is a framework-level classification,
not a physical observable.} The predicate D\_joint(A, B, Arch)
classifies whether a comparison architecture demands joint validity.
Operational sufficient conditions A--E (§6.3) connect observable
features of the setup (interference vs.~readout, comparison vs.~no
comparison) to D\_joint values in a reproducible way. But the mapping is
a VVV-QMRF classification of the experimental architecture, not a
measurement of a physical quantity.

\textbf{These constraints are typical of the measurement problem, not
unique to VVV-QMRF.} Every interpretation or extension of quantum
mechanics that addresses measurement introduces an unobservable
primitive:

{\def\LTcaptype{none} % do not increment counter
\begin{longtable}[]{@{}
  >{\raggedright\arraybackslash}p{(\linewidth - 4\tabcolsep) * \real{0.3333}}
  >{\raggedright\arraybackslash}p{(\linewidth - 4\tabcolsep) * \real{0.3333}}
  >{\raggedright\arraybackslash}p{(\linewidth - 4\tabcolsep) * \real{0.3333}}@{}}
\toprule\noalign{}
\begin{minipage}[b]{\linewidth}\raggedright
Framework
\end{minipage} & \begin{minipage}[b]{\linewidth}\raggedright
Central concept
\end{minipage} & \begin{minipage}[b]{\linewidth}\raggedright
Why not directly observable
\end{minipage} \\
\midrule\noalign{}
\endhead
\bottomrule\noalign{}
\endlastfoot
Copenhagen & Classical apparatus / Heisenberg cut & No formal definition
of what makes an apparatus classical \\
Many-Worlds (Everett) & Branching of worlds & No physical observable
distinguishes one branch \\
QBism & Agent's degree of belief & Subjective probability is not a
physical quantity \\
Relational QM (Rovelli) & Facts relative to observer & Relata are only
defined within the relation \\
Objective Collapse (GRW) & Collapse mechanism & Collapse events are
postulated, not yet detected \\
\textbf{VVV-QMRF} & \textbf{Self-certification / D\_joint} &
\textbf{K-side properties are not ρ-side observables (K ≠ H)} \\
\end{longtable}
}

What distinguishes VVV-QMRF is that the constraints are explicitly
named, traced to a single architectural commitment (K ≠ H), and given
operational bridges (ODC\_K, conditions A--E) that define what can and
cannot be tested under the current formalism.

\begin{center}\rule{0.5\linewidth}{0.5pt}\end{center}

\subsection{References}\label{references}

\begin{enumerate}
\def\labelenumi{\arabic{enumi}.}
\tightlist
\item
  von Neumann, J. (1932). \emph{Mathematical Foundations of Quantum
  Mechanics.} Princeton University Press.
\item
  Bohr, N. (1935). Can Quantum-Mechanical Description of Physical
  Reality Be Considered Complete? \emph{Physical Review} 48, 696--702.
\item
  Heisenberg, W. (1958). \emph{Physics and Philosophy.} Harper \& Row.
\item
  Wigner, E.P. (1961). Remarks on the mind-body question. In: \emph{The
  Scientist Speculates,} Good (ed.), pp.~284--302.
\item
  Rovelli, C. (1996). Relational quantum mechanics. \emph{International
  Journal of Theoretical Physics} 35, 1637--1678.
\item
  Prasad, H.S. (2023). Buddhist Pramāṇa-Epistemology. \emph{Studia
  Humana} 12(1-2), 21--52. DOI: 10.2478/sh-2023-0004.
\item
  Zurek, W.H. (2003). Decoherence, einselection, and the quantum origins
  of the classical. \emph{Reviews of Modern Physics} 75, 715--775.
\item
  Frauchiger, D. and Renner, R. (2018). Quantum theory cannot
  consistently describe the use of itself. \emph{Nature Communications}
  9, 3711.
\item
  Brukner, C. (2018). A no-go theorem for observer-independent facts.
  \emph{Entropy} 20(5), 350.
\item
  Proietti, M. et al.~(2019). Experimental test of local observer
  independence. \emph{Science Advances} 5(9), eaaw9832. DOI:
  10.1126/sciadv.aaw9832.
\item
  Bong, K.W. et al.~(2020). A strong no-go theorem on the Wigner's
  friend paradox. \emph{Nature Physics} 16, 1199--1205. DOI:
  10.1038/s41567-020-0990-x.
\item
  Jordan, A.N. and Siddiqi, I.A. (2024). \emph{Quantum Measurement
  Theory and Practice.} Cambridge University Press. DOI:
  10.1017/9781009103909.
\item
  Baumann, V. and Brukner, Č. (2024). Wigner's friend's memory and the
  no-signaling principle. \emph{Quantum} 8, 1481. arXiv:2305.15497.
\item
  VietVunVut (Viet - Nguyen Xuan) (2026). Has Every Wigner's Friend
  Experiment Been Blind to a Geometric Degree of Freedom? Submitted to
  arXiv 2026-05-27 (arXiv ID pending confirmation). Internal ref:
  \texttt{papers/paper\_002/manuscript.md}.
\item
  VietVunVut (Viet - Nguyen Xuan) (2026). VVV-QMRF Class C. Working
  Paper v2.0. DOI: 10.5281/zenodo.20289261. Zenodo.
\end{enumerate}

\begin{center}\rule{0.5\linewidth}{0.5pt}\end{center}

\subsection{Appendices}\label{appendices}

\textbf{Appendix A} --- K9\_E reproduction scripts:
\texttt{documents/research\_documents/project\_vvv\_qmrf\_class\_c/07\_fits/}
\textbf{Appendix B} --- K\_Space\_Axiomatization v2.4:
\texttt{documents/research\_documents/meta\_architecture/K\_Space\_Axiomatization.md}
\textbf{Appendix C} --- AHP audit footprint:
\texttt{documents/research\_documents/anti\_hallucinations/}

\begin{center}\rule{0.5\linewidth}{0.5pt}\end{center}

\subsection{Draft phase status}\label{draft-phase-status}

{\def\LTcaptype{none} % do not increment counter
\begin{longtable}[]{@{}
  >{\raggedright\arraybackslash}p{(\linewidth - 4\tabcolsep) * \real{0.3333}}
  >{\raggedright\arraybackslash}p{(\linewidth - 4\tabcolsep) * \real{0.3333}}
  >{\raggedright\arraybackslash}p{(\linewidth - 4\tabcolsep) * \real{0.3333}}@{}}
\toprule\noalign{}
\begin{minipage}[b]{\linewidth}\raggedright
Phase
\end{minipage} & \begin{minipage}[b]{\linewidth}\raggedright
Sections
\end{minipage} & \begin{minipage}[b]{\linewidth}\raggedright
Status
\end{minipage} \\
\midrule\noalign{}
\endhead
\bottomrule\noalign{}
\endlastfoot
\textbf{P1} & Skeleton + Abstract (8 paras) + §4 K-Space Architecture &
\textbf{COMPLETE 2026-05-28} \\
\textbf{P2} & §5 K9\_E Postulate + §7 Empirical Status &
\textbf{COMPLETE 2026-05-28} \\
\textbf{P3} & §8 K9-S12 + §10 T4-H + 3-OBS & \textbf{COMPLETE
2026-05-28} \\
\textbf{P4} & §9 BB+FR + §11 φ-map Class D & \textbf{COMPLETE
2026-05-28} \\
\textbf{P5} & Carry-over §1, §2, §3 (headers) + §6 + §12 + §13 &
\textbf{COMPLETE 2026-05-28} \\
\textbf{P6} & Abstract finalization + header update & \textbf{COMPLETE
2026-05-28} \\
\textbf{P7} & CHANGELOG + README + AHP footprint verified &
\textbf{COMPLETE 2026-05-28} \\
\end{longtable}
}

\begin{center}\rule{0.5\linewidth}{0.5pt}\end{center}

\emph{VVV-QMRF Working Paper v3.0 --- All phases P1-P7 complete
(2026-05-28). All mini-RCA PASS (4.57--4.6/5). Average 4.58/5.}
\emph{Promoted from draft to final 2026-05-28. Abstract (8 para) +
§1--§13 (\textasciitilde14,500w). Supporting files: CHANGELOG.md,
README.md, plan.} \emph{References {[}13{]} and {[}14{]} finalized;
{[}14{]} arXiv ID pending confirmation post-submission.} \emph{All
formal claims classified by evidence level. Class C = structurally
testable, empirically UNCONFIRMED. Class D = proposed.}

\end{document}
